% master.tex : master-fil for projektet
% ------------------------------------------------------------------------------
% Dette er hovedfilen for projektet, hvori indhold fra alle input-filer (tekst,
% billeder, litteraturdatabaser, osv.) samles

% Dokumenttypen 'book' er valgt pga. dens mange fleksible indstillinger
% Se https://tex.stackexchange.com/a/36989/118167
\documentclass[11pt,a4paper,twoside,openright,danish]{book}

% Variabler, som bruges til automatisk at indsætte titel, forfattere, osv. på
% forsiden og titelbladet.
\def \projecttitle       {Optimering}
\def \projectsubtitle    {Lineær programmering}
\def \projecttheme       {Lineær algebra}
\def \projectdegree      {Matematik}
\def \projectperiod      {Forårssemesteret 2020}
\def \projectnumber      {P2}
\def \projectgroup       {B303c}
\def \projectauthors     {
	Jens Koldkur\\
	Jonas Bach Christensen\\
	Julie Havbo Lund\\
	Mads Bjerregaard Kjær\\
	Mathias Gammelgaard\\
	Rasmus Jespersgaard
}
\def \projectsupervisors {
  Horia Cornean
}

% Preamblet indeholder alle de indstillinger og makroer, som skal indsættes for
% hovedindholdet, og i denne skabelon samles det i filen aaumath.sty, som
% definerer en pakke, der kan indlæses med \usepackage.
\usepackage{aaumath}

% Dokumentets indhold indsættes mellem \begin- og \end-makroerne for
% 'document'-blokken
\begin{document}

% Dokumentets 'front matter' tælles ikke med ifm. antal sider og nummereres med
% romerske tal. Herunder hører f.eks. forsiden, titelbladet, forordet og
% indholdsfortegnelsen.
\frontmatter
% incl/misc/frontpage.tex : rapportens forside
% ------------------------------------------------------------------------------


\backgroundsetup{
  scale = 1,
  angle=0,
  opacity=1,
  contents = {
    \includegraphics[width=\paperwidth,height=\paperheight]{fig/img/aau/waves.pdf}
  }
}
\BgThispage
\pdfbookmark[0]{Forside}{forside}
\begin{titlepage}
  \centering
  \phantom{}
  \vspace{2cm}

  % AAU-segl
  \begin{minipage}[c]{0.2\paperwidth}
    \centering
    \makebox[0pt]{
      % fig/tikz/aau-badge.tex : AAU-logo til forsiden
% ------------------------------------------------------------------------------

\begin{tikzpicture}
  % Tegn hvid cirkel og tilføj det gennemsigtige, blå logo ovenpå
  \node[circle,color=white,fill=white,minimum size=1.175\textwidth] at (0,0) {};
  \node at (0,0) {\includegraphics[width=\textwidth]{fig/img/aau/logo-circle.pdf}};
\end{tikzpicture}

    }
  \end{minipage}

  % Hovedindhold
  \vspace{4cm}
  {\fontfamily{bch}\selectfont
    \fboxsep0pt\colorbox{white}{
      \begin{minipage}{\textwidth}
        \centering
        \color{AAUblue1}

        \vspace{2em}
        {\Huge\bfseries\projecttitle}

        {\Large\bfseries\projectsubtitle}

        \bigskip
        \parbox{\textwidth}{\centering\large\projectauthors}

        \bigskip
        {\projectdegree}
        \vspace{2em}
      \end{minipage}
    }
  }

\end{titlepage}

% incl/misc/titlepage.tex : rapportens titelblad
% ------------------------------------------------------------------------------
% Titelbladet genereres af makroen \aautitlepage, som er defineret i
% /incl/pre/ext/aautitlepage.sty


\pdfbookmark[0]{Titelblad}{titelblad}
\aautitlepage{
  \projectinfo{
    \projecttitle
  }{
    \projecttheme
  }{
    \projectperiod
  }{
    Gruppe \projectgroup
  }{
    \parbox[t]{\textwidth}{\projectauthors}
  }{
    \parbox[t]{\textwidth}{\projectsupervisors}
  }{
    \today
  }
}{
  \textbf{Institut for Matematiske Fag}\\
  Skjernvej 4A\\
  DK-9220 Aalborg Ø\\
  \href{http://math.aau.dk}{http://math.aau.dk}
}{
  % incl/misc/abstract.tex : projektets abstract
% ------------------------------------------------------------------------------
% Et abstract er et kort resume af rapporten, som vises på titelbladet
This paper contains an examination on linear programming problems and their geometrical solutions in relation to the simplex method. 
The paper contains a theoretical introduction to linear algebra together with important definitions and theorems related to this field of study. 
Afterwards the solution of linear programming problems and notations are introduced, which leads into the geometrical representation of  linear programming problems through polyhedrons. Here it is proven that an optimal solutions exists in a corner of such a geometrical figure.
Lastly the report provides a link between the use of the simplex method, which is a solution algorithm for linear programming problems, and the geometrical representation, aswell as the practical application of this.
Different implementations of the simplex method are illustrated her aswell and discussed in regards to their time complexity. 
}

% incl/misc/contents.tex :
% ------------------------------------------------------------------------------

\pdfbookmark[0]{Indhold}{indhold}

% Indstillingerne i denne fil er grupperet, så de ikke påvirker andre dele
\begingroup

% Slå twoside fra midlertidigt for at undgår sideskift
\makeatletter
\@twosidefalse
\makeatother

% Placer indholdsfortegnelsen på sin egen side (bedst når den kun fylder en side)
\tableofcontents
\clearpage

% Placer oversigtslisterne på efterfølgende sider
%\let\clearpage\relax
%\listoffigures
%\listoftables
%\listofalgorithms
%\lstlistoflistings

\endgroup


% Dokumentets 'main matter' (hovedindhold) er der, hvor det meste indhold skal
% sættes ind. Sider og overskrifter er nummererede med arabiske tal.
\mainmatter

% Input-filer bør opdeles således, at hver fil svarer til et kapitel. Makroen
% \include indsætter et sideskift og indholdet fra den givne stil.

% INDLEDNING
	\chapter{Forord}

Jens Koldkur,
Jonas Bach Christensen,
Julie Havbo Lund,
Mads Bjerregaard Kjær,
Mathias Gammelgaard og
Rasmus Jespersgaard.
	
2. semester 
	
bacheloruddannelsen Matematik på Aalborg Universitet
	
2020
	
udarbejdet fra 3 februar til 27 maj 2020
	
kravende 
	
temaet 
	
Rapporten er nummereret i afsnit, definitioner, figurer m.m.
Beviser afsluttes med $\square$, som indikerer, at beviset er gennemført. \\\\
% \square skulle have været \qed!

En tak til vores vejleder Horia Cornean
	\chapter{Indledning}
%To til tre sider 
%Vi skal argumentere det er interessant at kigge på lineær programmering
%Kort beskrivelse af hvad vi har lavet 
%Eventuelt 
%   Kapitel 1: bla bla 
%   Kapitel 2: bla bla 
%   m.m.
Optimeringsproblemer er en fundamental del af matematikken, og i den anvendte matematik udgøre disse en vægtig del af matematikkeres virke.
I denne rapport vil liniære optimeringsproblemer undersøges. Dele af forskningsfeltet betragter disse som havende ophav ved Danzig og dennes udarbejdelse af simplexmetoden(ref) i 1947.
Der er dog tilstadighed diskussion om hvor liniære programeringsproblemer som matematisk felt stammer fra og andre peger på Sovjetunionen og behovet for at optimere produktionen i en planøkonomisk kontekst som ophav.(ref2)  
Forudsætningen for at Danzig arbejdede med liniære programeringsproblemer var hans ansættelse ved Royal Airforc(ref) og en applikation blev i forbindelse med den vestligt kontrolerede sektor af Berlin.
Vestberlin lå dybt inde i den sovjetisk kontrolerede sektor af Tyskland og som følge af spændingerne mellem vesteuropa og Sovjet besluttede Josef Stalin at lukke for forsyninger via jorden til Berlin.
Det blev derfor nødvendigt for vestmagterne at etablere en luftbro til Berlin der gjorde det muligt trods de sovjetiske restriktioner at indføre basale fornødehender til de $2,5$ milloner mennesker der levede i byen.
Da en sådan måde at indføre vare på er ressourceintensiv, blev der derfor behov for at gøre dette mest effektivt underlagt de logistiske restriktioner.
Det er derfor et matematisk felt der i sin brug er nært bundet op på praktiske problemstillinger dette værende økonomiske og logistiske.
Det er derfor også et felt der har været nært knyttet op til militære problemstillinger som Danzigs ansættelsesforhold også antyder.
Danzigs bidrag til forskningsfeltet, altså simplexmetoden, er dog stadigt en fundamental del af hvordan løsningen af liniære programmeringsproblemer behandles i dag.
Projektets primære arbejdsområde vil derfor omhandle simplex-metoden samt den geometriske fremstilling af liniære programeringsproblemer.
Det skal dog nævnes at der ligeledes findes optimeringsproblemer der ikke kan beskrives igennem liniære sammenhænge, en afgrænsning bliver her at disse ikke vil beskrives i dette indeværende projekt.
Problemformuleringen bliver derfor som følger:
Problemformulering:
Hvordan kan liniære porgrameringsproblemer anskues geometrisk og hvordan kan disse løses ved hjælp af simplexmetoden?
Med henblik på at besvare denne vil følgende emneområder bliver inddraget:
\begin{itemize}
\item En introduktion til liniær algebra: Afsnittet har til hensigt at indføre passende notation, samt introducere væsentlige metoder og begreber til løsning af liniære programeringsproblemer.
\item Liniære programeringsproblemer: Projektets primære emneområde. Det vil her introduceres hvordan disse fremstilles matematiske samt interessante egenskaber herved.
\item Geometrisk fremstilling: her skal der stå et eller andet når vi ved hvad der er med.
\item Simplexmetoden: Beskrivelse og bevis af Simplex som løsningsalgoritme.
\end{itemize}
	
% LINEÆR ALGEBRA 
	\chapter{Lineær algebra}
%
For at kunne behandle et optimeringsproblem, undersøges først en række elementære begreber indenfor lineær algebra, der er anvendelige i forbindelse med problemet.
Dette kapitel tager udgangspunkt i \citep[side 3-188 og 227-259]{lial}, hvis ikke andet er angivet.
%
\section{Matricer og vektorer}
%  
% Noget indledningsvist - hvorfor er matricer/vektorer vigtige? 
%
I forbindelse med løsning af lineære ligningssystemer og optimeringsproblemer er det nødvendigt at introducere \textbf{matricer} og tilhørende notation. 
%
\begin{defn}{}{}
En matrix er en rektangulær tabel, hvis elementer er \textbf{skalarer}, der er $k \in \R$. 
En matrix med $m$ rækker og $n$ søjler har størrelsen $m \times n$.
Hvis $m=n$, er matricen kvadratisk, noteret $n \times n$.
Elementet $a_{i,j}$ i $i$'te række og $j$'te søjle kaldes den $(i,j)$'te-indgang i matricen. 
\end{defn}
\noindent
%
På figur \ref{fig:matrix_gen_eks} ses en $m \times n$ matrix $A$.
%
\begin{figure}[H]
	\begin{center}
$$
A=
\begin{bmatrix}
a_{1,1} & a_{1,2} & \ldots & a_{1,n} \\
a_{2,1} & a_{2,2} & \ldots & a_{2,n} \\
\vdots  & \vdots  & \ddots & \vdots \\
a_{m,1} & a_{m,2} & \ldots & a_{m,n} \\
\end{bmatrix}
$$
	\end{center}
	\caption{En $m \times n$ matrix $A$.}
	\label{fig:matrix_gen_eks}
\end{figure}
\noindent
%
Enhver række $i$ eller søjle $j$ kan betragtes som en \textbf{vektor}. 
En matrix med én række kaldes en \textbf{rækkevektor}, og en matrix med én søjle kaldes en \textbf{søjlevektor}. 
Med mindre andet er angivet opereres der med søjlevektorer.
I denne rapport noteres vektorer med små, fede bogstaver, såsom $\textbf{a}$ og $\textbf{b}$. 
Søjlevektorer noteres ofte med et indeks tilsvarende søjlen i matricen, eksempelvis:
%
$$
\textbf{a}_j= 
\begin{bmatrix}
a_{1,j} \\
a_{2,j} \\
\vdots \\
a_{i,j} \\
\end{bmatrix}.
$$ 
	% Matricer og vektorer
	\subsection{Specielle matricer} %skal måske hedde noget andet? 

Den ene af to specielle matricer, der benyttes i dette projekt, er \textit{nulmatricen}. I en nulmatrice er alle elementer $0$, og en $m \times n$ nulmatrice noteres $0_{m,n}$. 
%Er det korrekt notation? Fandt det på wiki hehehehe

Den anden specielle matrice er \textit{identitetsmatricen}. 

\begin{defn}{}{}
En $n \times n$, $n \in \Z^+$, identitetsmatrice $I_n$ består af alle standardvektorer $\textbf{e}_1, textbf{e}_2, \ldots, textbf{e}_n$ i $R^n$, således 

$
I_n=
\begin{bmatrix}
\textbf{e}_1 & textbf{e}_2 & \ldots & textbf{e}_n
\end{bmatrix}
$. 

\end{defn}

Standardvektorerne i $\R^n$ er defineret ved 
%Kan evt. kommer inden definitionen, men kan bedst lide denne rækkefølge. 
$$
\textbf{e}_1=
\begin{bmatrix}
1 \\ 
0 \\ 
\vdots \\
0
\end{bmatrix}
\text{, }
\textbf{e}_2=
\begin{bmatrix}
0 \\ 
1 \\ 
\vdots \\
0
\end{bmatrix}
\text{, }
\ldots
\text{, }
\textbf{e}_n=
\begin{bmatrix}
0 \\ 
0 \\ 
\vdots \\
1
\end{bmatrix}
\text{. }
$$


	\subsection{Matrixsum og skalering}
% 
Der er behov for at lave opperationer på en matrix i forbindelse med løsning af lineære programmeringsproblemer. 
To af disse er \textit{matrixsum} og \textit{skalering}.
%
\begin{defn}{}{mxsum} % Jeg synes substration skal være fed
Lad $A$ og $B$ være $m \times n$ matricer.
\textbf{Summen} af $A$ og $B$ er $m \times n$ matricen $A + B$, hvor den $(i,j)$'te indgang er $a_{i,j} + b_{i,j}$.
Ligeledes er subtraktion mellem matricerne $A$ og $B$ muligt, og indgangene i den resulterende matrix er $a_{i,j} - b_{i,j}$.
\\
Lad nu $c$ være en skalar.
\textbf{Skalering}, også kaldet skalarmultiplikation, af $A$, noteret $cA$, er en $m \times n$ matrix, hvor den $(i,j)$'te indgang er $ca_{i,j}$.
\end{defn}
\noindent
%
Bemærk, at $1A = A$, $-1A = -A$ og $0A = O$.
\\
%
\begin{eks}
Lad 
\begin{align*}
A= 
\begin{bmatrix}
3	&	-3	&	1\\
2	&	0	&	4
\end{bmatrix}
\text{\phantom{---}og\phantom{---}}
B= 
\begin{bmatrix}
-1	&	4	&	1\\
5	&	3	&	2
\end{bmatrix}.
\end{align*}
Jævnfør \ref{defn:mxsum} er
\begin{align*}
3A= 
\begin{bmatrix}
9	&	-9	&	3\\
6	&	0	&	12
\end{bmatrix}
\text{,\phantom{---}}
-B= 
\begin{bmatrix}
1	&	-4	&	-1\\
-5	&	-3	&	-2
\end{bmatrix}
\end{align*}
og
\begin{align*}
3A+B= 
\begin{bmatrix}
9	&	-9	&	3\\
6	&	0	&	12
\end{bmatrix}
+ 
\begin{bmatrix}
-1	&	4	&	1\\
5	&	3	&	2
\end{bmatrix}
=
\begin{bmatrix}
8	&	-5	&	4\\
11	&	3	&	14
\end{bmatrix}.
\end{align*}
\end{eks}
%
Ud fra disse operationer har matricer en række egenskaber, som ses i \ref{thm:mxprop}.
%Nedenstående definition er ordret den samme som i bogen (s. 6), men vi kan jo nærmest ikke skrive den og lignende sætninger anderledes; skal vi ikke bare lige vende det med Horia, tror I? 
\begin{thm}{}{mxprop}
Lad $A$, $B$ og $C$ være $m \times n$ matricer, og lad $c$ og $k$ være skalarer.
Så har matricer følgende egenskaber:
\begin{enumerate}[label=(\alph*)]
\item $A + B = B + A$.
\item $(A + B) + C = A + (B + C)$.
\item $A + O = A$.
\item $A + (-A) = O$.
\item $(ck)A = k(cA)$.
\item $c(A + B) = cA + cB$.
\item $(c + k)A = cA + kA$.
\end{enumerate}
\end{thm}
%
\begin{proof} % Skal dette bevis være i punktform? det skal det jo vel egenlig 
Lad $A$, $B$ og $C$ være $m \times n$ matricer, og lad $c$ og $k$ være skalarer.
Betragt $A + B$ og $B + A$ og bemærk jævnfør \ref{defn:mxsum}, at de $(i,j)$'te indgange bliver henholdsvis $a_{i,j} + b_{i,j}$ og $b_{i,j} + a_{i,j}$, hvilket beviser (a).
Ligeledes ses det, at $(a_{i,j} + b_{i,j}) + c_{i,j} = a_{i,j} + (b_{i,j} + c_{i,j})$, hvilket beviser (b).
Endvidere kan $0$ isoleres i $a_{i,j} + 0 = a_{i,j}$, således $a_{i,j} + (-a_{i,j}) = 0$, hvilket beviser (c) og (d), da alle indgange i en nulmatrix er $0$.
\\\\
%
Betragt skaleringen af $A$ med skalaren $ck$, som er et produkt af skalarene $c$ og $k$, og skaleringen af $A$ med $c$ og herefter $k$.
For indgangene i disse skaleringer findes ligheden $(ck)a_{i,j} = k(ca_{i,j})$, hvilket beviser (e).
Betragt ligeledes (f) og bemærk, at $c(a_{i,j} + b_{i,j}) = ca_{i,j} + cb_{i,j}$, hvilket beviser (f).
Det samme gøres for ligheden i (g), hvor ligheden $(c + k)a_{i,j} = ca_{i,j} + ka_{i,j}$ opstilles, hvilket beviser (g).
\end{proof}
\noindent
\\
%
Da vektorer er matricer med enten én række eller én søjle er samme operationer mulige og vektorer besidder derfor de samme egenskaber.
%
% Mere sssssssssssssssssssssssssssssssssssss
\\\\
\phantom{Hejmeddig}
\\\\
	\subsection{Matrix transponering}
\begin{defn}{}{mxtrans}
Lad $A$ være en $m \times n$ matrix. \textbf{Transponeringen} til $A$, er en $n \times m$ matrix, noteret $A^T$, hvor $A$'s $(i,j)$te indgang er den $(j,i)$te indgang i $A^T$.
\end{defn}
\begin{eks}
Betragt nu $A^T$ og $(A+B)^T$, hvor 
\begin{align*}
A= 
\begin{bmatrix}
3	&	-3	&	1\\
2	&	0	&	4
\end{bmatrix}
\text{\phantom{---}og\phantom{---}}
B= 
\begin{bmatrix}
-1	&	4	&	1\\
5	&	3	&	2
\end{bmatrix}.
\end{align*}

\end{eks}
	\subsection{Linearkombination}
%
\textbf{Linearkombination} er en centralt del af lineær algebra, da det benyttes i forbindelse med løsning af ligningssystemer. 
%
\begin{defn}{}{mxlinkom}
Givet vektorer $\mathbf{u}_1, \mathbf{u}_2, \ldots, \mathbf{u}_k$ med størrelsen $m \times n$, så er linearkombinationen af $\mathbf{a}_1, \mathbf{u}_2, \ldots, \mathbf{u}_k$ en vektor $\mathbf{v}$ på formen 
%
\begin{align*}
c_1\mathbf{u}_1+c_2\mathbf{u}_2+\ldots+c_k\mathbf{u}_k=\sum\limits_{i=1}^k c_i\mathbf{u}_i=\mathbf{v},
\end{align*}
%
med størrelsen $m \times n$, hvor $c_1, c_2, \ldots, c_k$ er skalarerne.
Skalarerne kaldes linearkombinationens koefficienter.
\end{defn}
%
% KILDE primært: https://www.statlect.com/matrix-algebra/linear-combinations
%
\noindent 
%
En linearkombination af én vektor er blot en skalering af vektoren.
\\
%
%%%%%%%%%%%%%%%%%%%%%%%%%%%%%%%%%%%%%%%%%%%
%
\begin{eks}
Lad 
$$
\textbf{a}=
\begin{bmatrix}
5 \\
9 \\
4
\end{bmatrix}
\text{, }
\textbf{b}=
\begin{bmatrix}
8  \\
12 \\ 
3
\end{bmatrix}
\text{og }
\textbf{c}=
\begin{bmatrix}
0 \\
0 \\
7
\end{bmatrix},
$$
samt $c_a=-4$, $c_b=100$ og $c_c=2$. 
Jævnfør definition \ref{defn:mxlinkom} er linearkombinationen af $\textbf{a}$, $\textbf{b}$ og $\textbf{c}$:
$$
(-4)
\begin{bmatrix}
5 \\ 
9 \\ 
4
\end{bmatrix}
+
100
\begin{bmatrix}
8 \\
12 \\ 
3
\end{bmatrix}
+
2
\begin{bmatrix}
0 \\
0 \\ 
7
\end{bmatrix}
=
\begin{bmatrix}
780 \\
1164 \\ 
298
\end{bmatrix}
\text{.}
$$
\end{eks}
%
%%%%%%%%%%%%%%%%%%%%%%%%%%%%%%%%%%%%%%%%%%%%%%
%
Enhver vektor 
\begin{align*}
\textbf{a}=
\begin{bmatrix}
a_1 \\
a_2 \\
\vdots \\
a_n
\end{bmatrix}
\end{align*}
%
i $\R^n$ kan skrives som linearkombination af $I_n$ i $\R^n$. 
Givet to ikke-parallelle vektorer $\mathbf{u}$ og $\mathbf{v}$ i $\R^2$, så er enhver anden vektor i $\R^2$ en linearkombination af $\mathbf{u}$ og $\mathbf{v}$. 
Hvis koefficienterne i en linearkombination er vektorer, så er der tale om et \textbf{matrix-vektorprodukt}.
%
%
%
%
\subsection{Matrix-vektorprodukt}
%
Som udvidelse til linearkombinationen kan der tales om et\textbf{matrix-vektorprodukt}.
%
%Altså, jeg vil måske mene, at nedenstående definition minder for meget om den i bogen (s. 19). Men igen... "Det er heller ikke ham, der har fundet på det." - Mads Vejleder, 2019. 
%
\begin{defn}{}{mvp}
Lad $A$ være en $m \times n$ matrix og \textbf{u} være en $n \times 1$ vektor. 
Matrix-vektorproduktet af $A$ og $\textbf{u}$ defineres som linearkombinationen af søjlerne i $A$, hvis koefficienter er de tilsvarende komponenter af $\textbf{u}$, hvilket noteres $A\textbf{u}$. 
Dermed er
\begin{align*}
A\textbf{u} =u_1\textbf{a}_1 + u_2\textbf{a}_2 + \cdots + u_n\textbf{a}_n.
\end{align*}
\end{defn}
\noindent
For at $A\textbf{u}$ eksisterer, skal der være lige mange komponenter i \textbf{v}, som der er søjler i $A$. 
Et eksempel på et matrix-vektorprodukt mellem en matrix $A$ og en vektor \textbf{u} ses i eksempel \ref{Matrix-vektor}.
\\
%
\begin{eks}
\label{Matrix-vektor}
%
Lad
$$A=
\begin{blockarray}{c c c c}
\begin{block}{[c c c c]}
2 & 5 & 6 \\
4 & 7 & 3\\
1 & 4 & 1\\
2 & 3 & 8\\
\end{block}
\end{blockarray}
%
\text{ og }
%
\textbf{u}=
\begin{bmatrix}
4 \\
2 \\
3 \\ 
\end{bmatrix}.
$$
%
Eftersom $A$ har 3 søjler og \textbf{u} har 3 komponenter, så eksisterer $A\textbf{u}$ og findes ved
$$
A\textbf{u}=
\begin{bmatrix}
2 & 5 & 6 \\
4 & 7 & 3\\
1 & 4 & 1\\
2 & 3 & 8\\
\end{bmatrix}
\begin{bmatrix}
4 \\
2 \\
3 \\ 
\end{bmatrix}
=4
\begin{bmatrix}
2\\
4\\
1\\
2\\
\end{bmatrix}
+2
\begin{bmatrix}
5\\
7\\
4\\
3\\
\end{bmatrix}
+3
\begin{bmatrix}
6\\
3\\
1\\
8\\
\end{bmatrix}
=
\begin{bmatrix}
8\\
16\\
4\\
8\\
\end{bmatrix}
+
\begin{bmatrix}
10\\
14\\
8\\
6\\
\end{bmatrix}
+
\begin{bmatrix}
18\\
9\\
3\\
24\\
\end{bmatrix}
=
\begin{bmatrix}
36\\
39\\
15\\
38\\
\end{bmatrix}.
$$
\end{eks}
	\subsection{Matrixmultiplikation}

Det er ofte nødvendigt at multiplicere matricer med andre matricer. 
\begin{defn}{}{mxmulti}
Lad $A$ være en $m \times n$ matrix og $B$ være en $n \times p$ matrix.
\textbf{Matrixproduktet} $AB$  er en $m \times p$ matrix 
$$
C=
\begin{bmatrix}
A\textbf{b}_1 & A\textbf{b}_2 & \ldots & A\textbf{b}_p
\end{bmatrix}\text{,}
$$
hvor $A\textbf{b}_j$ er den $j$'te søjle i $C$.
\end{defn}
\noindent
%
Som det fremgår af definition \ref{defn:mxmulti}, skal antallet af søjler i $A$ tilsvare antallet af rækker i $B$, for at $AB$ er defineret og matrixmultiplikation dermed er muligt. 
Bemærk, at selvom $AB$ er defineret, er $BA$ ikke nødvendigvis defineret. 
Hvis både $AB$ og $BA$ er defineret, er de desuden ikke nødvendigvis lig hinanden. 
For at beregne værdien af komponenterne i matrixproduktet $AB$ kan følgende sammenhæng mellem den $i$'te række i $A$ og den $j$'te søjle i $B$ betragtes: 
\begin{align*}
a_{i1}b_{1j} + a_{i2}b_{2j} + \cdots + a_{in}b_{nj}
\text{.}
\end{align*}
%\\
%
\begin{eks}
Lad 
$$
A=
\begin{bmatrix}
7 & 9 & 1 \\
4 & 0 & 0 \\
8 & 5 & 5
\end{bmatrix}
\text{og }
B=
\begin{bmatrix}
2 & 3 \\
5 & 1 \\
0 & 1 
\end{bmatrix}
\text{.}
$$
Ved at betragte $A$ og $B$ vides det, $AB$ er defineret og er en $3 \times 2$ matrix. 
Først findes søjlerne $A\textbf{b}_1$ og $A\textbf{b}_2$:
$$
A\textbf{b}_1=
\begin{bmatrix}
7 & 9 & 1 \\
4 & 0 & 0 \\
8 & 5 & 5
\end{bmatrix}
\begin{bmatrix}
2 \\
5 \\
0
\end{bmatrix}
=
\begin{bmatrix}
14 + 45 + 0 \\
8 + 0 + 0 \\
16 + 25 + 0
\end{bmatrix}
=
\begin{bmatrix}
59 \\
8 \\
41
\end{bmatrix}
$$
og
$$
A\textbf{b}_1=
\begin{bmatrix}
7 & 9 & 1 \\
4 & 0 & 0 \\
8 & 5 & 5
\end{bmatrix}
\begin{bmatrix}
3 \\
1 \\
1
\end{bmatrix}
=
\begin{bmatrix}
21 + 9 + 1 \\
12 + 0 + 0 \\
24 + 5 + 5
\end{bmatrix}
=
\begin{bmatrix}
31 \\
12 \\
34
\end{bmatrix}
\text{.}
$$
Dermed er 
$$
AB=
\begin{bmatrix}
A\textbf{b}_1 & A\textbf{b}_2
\end{bmatrix}
=
\begin{bmatrix}
59 & 31 \\
8 & 12 \\
41 & 34
\end{bmatrix}.
$$
\end{eks}
Da matricen $B$ inddeles i sine søjlevektorer, er grundprincippet i matrixmultiplikation at finde matrix-vektorprodukter.\\\\
En række egenskaber gælder for matrixmultiplikation, hvilket er summeret i sætning \ref{thm:mxmulti}.

\begin{thm}{}{mxmulti}
Lad $A$ og $B$ være $k \times m$ matricer, $C$ være en $m \times n$ matrix, og $P$ og $Q$ være $n \times p$ matricer. Følgende udsagn er dermed sande:
\begin{enumerate}[label=(\alph*)]
\item $s(AC)=(sA)C=A(sC)$ for en vilkårlig skalar $s$.
\item $A(CP)=(AC)P$.
\item $(A+B)C=AC+BC$.
\item $C(P+Q)=CP+CQ$.
\item $I_kA=A=AI_m$.
\item Produktet af en nulmatrix og en vilkårlig matrix er altid $O$.
\item $(AC)^T=C^TA^T$.
\end{enumerate}
\end{thm}
\begin{proof}
Lad $A$ og $B$ være $k \times m$ matricer, $C$ være en $m \times n$ matrix, og $P$ og $Q$ være $n \times p$ matricer. 
\\\\
(a) Den $(i,j)$'te indgang i matrixprodukterne $s(AC)$ og $(sA)C$ har formen $sa_{i,1}c_{1,j} + sa_{i,2}c_{2,j} + \cdots + sa_{i,m}c_{m,j}$, hvilket er det samme som $a_{i,1}sc_{1,j} + a_{i,2}sc_{2,j} + \cdots + a_{i,m}sc_{m,j}$, der er en indgang i $A(sC)$. 
Da indgangene er ens, er matrixproduktet det samme i de tre tilfælde.
\\\\
(b) Bemærk, at $A$ ikke kan multipliceres med $P$, mens $C$ kan multipliceres med begge. 
Derfor kan matrixprodukterne $AC$ og $CP$ betragtes som to matricer, således at både $A(CP)$ og $(AC)P$ er $k \times p$ matricer. 
At komponenten i den $(i,j)$'te indgang i $A(CP)$ er ækvivalent med komponenten i den $(i,j)$'te indgang $(AC)P$ fås ved samme fremgangsmåde som vist i (a). 
% Måske er (b) ikke helt nok? 
\\\\
(c) Lad $\textbf{c}_j$ være en vilkårlig søjle i $C$. 
Det følger deraf, at $(A+B)\textbf{c}_j=A\textbf{c}_j+B\textbf{c}_j$ jævnfør sætning \ref{thm:mxvpro}. 
\\\\
(d) Lad $\textbf{p}_j$ og $\textbf{q}_j$ være søjler i henholdvis $P$ og $Q$. Det følger deraf, at $C(\textbf{p}_j+\textbf{q}_j)=C\textbf{p}_j+C\textbf{q}_j$ jævnfør sætning \ref{thm:mxvpro}. 
\\\\
(e) Da $I_k$ og $I_m$ henholdsvis korresponderer til antallet af rækker og søjler i $A$, er matrixprodukterne $I_kA$ og $AI_m$ defineret. 
Da $I_k\textbf{a}_j=\textbf{a}_j$ og $I_m\textbf{a}_i=I_m\textbf{a}_i$, er $I_kA=A$ og $AI_m=A$, hvormed $I_kA=AI_m$.
\\\\
(f) Eftersom enhver komponent i en vilkårlig matrix multipliceret med $0$ er $0$, følger udsagnet trivielt. 
\\\\
(g) Den $(i,j)$'te indgang i $(AC)^T$ er den $(j,i)$'te indgang i $AC$. De to udtryks ækvivalens vises ved sammen fremgangsmåde som i (a).
\end{proof}
	
	
	% Lineære ligningssystemer
	\section{Lineære ligningssystemer }
%
En \textit{lineær ligning} med variablerne $x_1,x_2,\ldots ,x_n$ er en ligning, som har formen
$$ c_1x_1+c_2x_2+\cdots+c_nx_n=b \text{,}$$ 
hvor $\{c_1,c_2,\ldots,c_n\} \in \R$. 
%
\begin{defn}{}{linlignsys}
Et \textbf{lineært ligningssystem} er en mængde af $m$ lineære ligninger med de samme $n$ variable, hvor $\{m, n\} \in \Z^+$.
Dette kan skrives på formen
\begin{align*}
c_{11}x_1+c_{12}x_2+&\cdots+c_{1n}x_n=b_1\\
c_{21}x_1+c_{22}x_2+&\cdots+c_{2n}x_n=b_2\\
&\vdots\\
c_{m1}x_1+c_{m2}x_2+&\cdots +c_{mn}x_n=b_m
\text{,}
\end{align*}
hvor $c_{ij}$ betegner koefficienten af $x_j$ i ligning $i$.
\end{defn}
\noindent
Løsningen på et sådant system er en vektor
\begin{center}
$
\mathbf{s} = 
\begin{bmatrix}
s_1\\
s_2\\
\vdots\\
s_n
\end{bmatrix},
$
\end{center}
i $\R ^n$, således at alle ligninger er opfyldt, når $x_i$ erstattes med $s_i$. Et eksempel på et lineært ligningssystem ses i \ref{eks:lignsys}.
%
\\
%
\begin{eks}
\label{eks:lignsys}
Understående er et lineært ligningssystem, bestående af tre lineære ligninger.
%
\begin{align*}
x_1-x_3-2x_4-8x_5&=-3 \\
-2x_1+x_3+2x_4+9x_5&=5 \\
3x_1-2x_3-3x_4-15x_5&=-9.
\end{align*}
%
% Løsningen på dette system ses i \ref{eks_gauss}.
%
%%%%%%%%%%%%%%%%%%%%%%%%%%%%%%%%%%%%%%%
%
\end{eks}
%
Alle lineære ligningssystemer har enten ingen, én eller uendeligt mange løsninger.
Ligninger i $\R^2$ kan repræsenteres ved linjer, og ligninger i $\R^3$ kan repræsenteres ved planer.
% Er det en eller et plan? Hess siger altid "en"... Hmmmmm.. 
Hvis disse ligninger er parallelle og forskellige, har systemet ingen løsning.
Hvis ligningerne er forskellige, så har systemet præcis én løsning, hvor ligningerne krydser hinanden. 
Hvis ligningerne er ens, har systemet uendeligt mange løsninger.
\\\\
%
Et system der har én eller flere løsninger kaldes \textit{konsistent}; ellers kaldes det \textit{inkonsistent}, da der ingen løsninger eksisterer.
%
%%%%%%%%%%%%%%%%%%%%%%%%%%%%%%%%%%%%%%%%%%%%%%%%%%%%%%%%%%%%%%%%%%%%%%%%%%%
%
\subsection{Ligningssystemer og matricer}
Et lineært ligningssystem kan opstilles som en \textit{matrixligning} $A\textbf{x}=\textbf{b}$, hvor
$$A=
\begin{bmatrix}
a_{1,1} & a_{1,2} & \cdots & a_{1,n}\\
a_{2,1} & a_{2,2} & \cdots & a_{2,n}\\
\vdots & \vdots & \ddots & \vdots\\
a_{m,1} & a_{m,2} & \cdots & a_{m,n}
\end{bmatrix}
\text{, } 
\textbf{x}=
\begin{bmatrix}
x_1\\
x_2\\
\vdots\\
x_n\\
\end{bmatrix}
\text{og }
\textbf{b}=\begin{bmatrix}
b_1\\
b_2\\
\vdots\\
b_m
\end{bmatrix}.
$$
%
A kaldes for \textit{koefficientmatricen}. 
%efter hårdt sammenleje bliver A matricen og b til totalmatricen
\textit{Totalmatrixen} $[A \mid \mathbf{b}]$ dannes ved at kombinere koefficientmatricen med vektoren \textbf{b}.
Totalmatrixen indeholder alle informationer, der skal bruges til at løse et lineært ligningssystem.
En matrixligning kan også opskrives ved
%
\begin{align*}
\sum^n_{i=1}{A_ix_i}=\textbf{b}.
\end{align*}
%
Totalmatricen opskrives på formen:
%
\begin{equation*}
[A \mid \mathbf{b}]=
\begin{blockarray}{cccccc}
x_1 & x_2 & \cdots & x_n & b \\
\begin{block}{[cccc|c]c}
a_{1,1} & a_{1,2} & \cdots & a_{1,n} & b_1 \\
a_{2,1} & a_{2,2} & \cdots & a_{2,n} & b_2 \\
\vdots & \vdots & \ddots & \vdots & \vdots \\
a_{m,1} & a_{m,2} & \cdots & a_{m,n} & b_{m}\\
\end{block}
\end{blockarray}.
\end{equation*}
%
\\
Ligningssystemet fra \ref{eks:lignsys} opskrives som totalmatricen i \ref{eks:lig_mat}. 
%
%
\begin{eks}\label{eks:lig_mat}
Med udgangspunkt i ligningssystemet fra \ref{eks:lignsys} følger det at
$$A=
\begin{bmatrix}
1 & 0 & -1 & -2 & -8\\
-2 & 0 & 1 & 2 & 9\\
3 & 0 & -2 & -3 & -15
\end{bmatrix}
\text{, } 
\textbf{x}=
\begin{bmatrix}
x_1\\
x_2\\
x_3\\
x_4\\
x_5
\end{bmatrix}
\text{, }
\textbf{b}=\begin{bmatrix}
-3\\
5\\
-9
\end{bmatrix}
$$
%
og totalmatricen 
%
\begin{equation*}
  [A \text{    } | \mathbf{b}] =
\begin{blockarray}{ccccccc}
x_1 & x_2 & x_3 & x_4 & x_5 & b \\
\begin{block}{[ccccc|c]c}
  1 & 0 & -1 & -2 & -8 & -3 \\
  -2 & 0 & 1 & 2 & 9 & 5 \\
  3 & 0 & -2 & -3 & -15 & -9 \\
\end{block}
\end{blockarray}.
\end{equation*}
%
\end{eks}

	%Matrixvektorproduct
\begin{defn}{}{}
\phantom{gdfs}\\Lad $A$ være en $m$ $x$ $n$ matrix og \textbf{v} være en $n$ $x$ $1$ vektor. \textbf{Matrix-vektor-produktet} af A og \textbf{v} defineres som den lineære kombination af søjlerne i A hvis koefficienter er de tilsvarende komponenter af \textbf{v}, og noteres A\textbf{v}. Det vil sige
\begin{align*}
A\textbf{v} =v_1\textbf{a}_1 + v_2\textbf{a}_2 + \cdots + v_n\textbf{a}_n
\end{align*}
\end{defn}
\noindent
For at $A\textbf{v}$ eksisterer, skal der være lige mange komponenter i \textbf{v} som søjler i $A$. Et eksempel på et matrix-vektor-produkt mellem en matrix $B$ og en vektor \textbf{u} ses i eksempel \ref{Matrix-vektor}.
\begin{eks}\label{Matrix-vektor}
\begin{center} 
$$B=
\begin{blockarray}{c c c c}
\begin{block}{[c c c c]}
2 & 5 & 6 \\
4 & 7 & 3\\
1 & 4 & 1\\
2 & 3 & 8\\
\end{block}
\end{blockarray}
%
\text{ og }
%
\textbf{u}=
\begin{bmatrix}
4 \\
2 \\
3 \\ 
\end{bmatrix}
$$
Eftersom $B$ har 3 søjler og \textbf{u} har 3 komponenter, så eksisterer $B\textbf{u}$ og findes ved
$$
B\textbf{u}=
\begin{bmatrix}
2 & 5 & 6 \\
4 & 7 & 3\\
1 & 4 & 1\\
2 & 3 & 8\\
\end{bmatrix}
\begin{bmatrix}
4 \\
2 \\
3 \\ 
\end{bmatrix}
=4
\begin{bmatrix}
2\\
4\\
1\\
2\\
\end{bmatrix}
+2
\begin{bmatrix}
5\\
7\\
4\\
3\\
\end{bmatrix}
+3
\begin{bmatrix}
6\\
3\\
1\\
8\\
\end{bmatrix}
=
\begin{bmatrix}
8\\
16\\
4\\
8\\
\end{bmatrix}
+
\begin{bmatrix}
10\\
14\\
8\\
6\\
\end{bmatrix}
+
\begin{bmatrix}
18\\
9\\
3\\
24\\
\end{bmatrix}
=
\begin{bmatrix}
36\\
39\\
15\\
38\\
\end{bmatrix}
$$
\end{center}
\end{eks}
\begin{eks}
test for noget
\end{eks}	
	% Løsning til lineære ligningssystemer
	\section{Løsning til lineære ligningssystemer}
%
Disse ligningssystemer løses ved hjælp af forskellige operationer, og i det følgende vil det introduceres, hvordan dette udføres.
%
\subsection{Elementære rækkeoperationer}
Der eksisterer tre operationer, som kan udføres på rækker i et lineært ligningssystem.
%
\begin{defn}{}{element}
Ved \textbf{ombytning} bytter to rækker plads i ligningssystemet, hvilket noteres
\begin{align*}
A \xrightarrow{R_i \leftrightarrow R_j} B, 
\end{align*}
hvor $A$ er matricen før rækkeoperationen udføres, og $B$ er matricen efter rækkeoperationen er udført. 
$R_j$ og $R_i$ er rækker i $A$.\\\\
Ved \textbf{skalering} multipliceres en række med et tal $c$, hvor $c \neq 0$, hvilket noteres
\begin{align*}
A \xrightarrow{R_i \rightarrow cR_i} B.
\end{align*}
Ved \textbf{udskiftning} udskiftes en række med rækken selv plus en skalering af en anden række, hvilket noteres
\begin{align*}
A \xrightarrow{R_i \rightarrow R_i + cR_h} B.
\end{align*}
%
\end{defn}
\noindent
Rækkeoperationerne kan udføres, når ligningssystemet er skrevet som en matrix. Eksempler på de elementære rækkeoperationer ses i \ref{eks1}.
\\
%
% Hvis vi bruger tilde skal den kommenteres hær
%
\begin{eks}\label{eks1}
Først laves ombytning mellem $R_1$ og $R_3$ i matricen. 
\begin{align*}
\begin{blockarray}{ccc}
\begin{block}{[cc|c]}
5 & 5 & 3 \\
2 & 2 & 1\\
3 & 4 & 8\\
\end{block}
\end{blockarray}
&\xrightarrow{R_1 \leftrightarrow R_3}
\begin{blockarray}{ccc}
\begin{block}{[cc|c]}
3 & 4 & 8\\
2 & 2 & 1\\
5 & 5 & 3\\
\end{block}
\end{blockarray}
\end{align*}
%
Derefter skaleres $R_2$ med $-2$. 
%
\begin{align*}
\begin{blockarray}{ccc}
\begin{block}{[cc|c]}
3 & 4 & 8\\
2 & 2 & 1\\
5 & 5 & 3\\
\end{block}
\end{blockarray}
&\xrightarrow{R_2 \rightarrow (-2)R_2}
\begin{blockarray}{ccc}
\begin{block}{[cc|c]}
3 & 4 & 8\\
-4 & -4 & -2\\
5 & 5 & 3\\
\end{block}
\end{blockarray}
\end{align*}
%
Slutteligt laves en udskiftning på $R_3$, hvor skaleringen $3R_1$ adderes med $R_3$. 
%
\begin{align*}
\begin{blockarray}{ccc}
\begin{block}{[cc|c]}
3 & 4 & 8\\
-4 & -4 & -2\\
5 & 5 & 3\\
\end{block}
\end{blockarray}
&\xrightarrow{R_3 \rightarrow 3R_1+R_3}
\begin{blockarray}{ccc}
\begin{block}{[cc|c]}
3 & 4 & 8\\
-4 & -4 & -2\\
14 & 17 & 27\\
\end{block}
\end{blockarray}
\end{align*}
\end{eks}
%
% Jeg kunne personligt godt tænke mig mere "matematik" i disse subsections, men jeg ved jo godt, at det er smag og behag. 
%
\subsection{Trappeform og reduceret trappeform}
En række i en matrix kaldes en \textit{nulrække}, hvis alle indgange er nul, og en \textit{ikke-nulrække}, hvis mindst én indgang i rækken er forskellig fra nul.
Dette bruges til definitionen af \textit{trappeform}.
%
\begin{defn}{}{}
En matrix er på \textbf{trappeform}, hvis den opfylder følgende kriterier:
\begin{enumerate}[label=(\alph*)]
\item Enhver \textbf{ikke-nulrække} ligger over alle nul rækker.
\item Den første ikke-nulindgang i en ikke-nulrække ligger i en søjle til højre for første ikke-nulindgange i forrige række.
\item Hvis en søjle har den første ikke-nulindgang i en række, så er alle indgangene i søjlen under den $0$.
\end{enumerate}
\end{defn}
%
\begin{eks}\label{eks:trappe}
Betragt $A$ og $B$. $A$ er på trappeform, mens $B$ ikke er på trappeform. 
Hvis der foretages en ombytning på række to og tre, så ville matrix $B$ også være på trappeform
%
\begin{align*}
A=
\begin{blockarray}{cccc}
\begin{block}{[ccc|c]}
2 & 4 & 8 & 2\\
0 & 5 & -2 & 5\\
0 & 0 & 2 & 7\\
0 & 0 & 0 & 0\\
\end{block}
\end{blockarray}
\text{ }
B=
\begin{blockarray}{cccc}
\begin{block}{[ccc|c]}
2 & 4 & 8 & 2\\
0 & 0 & -2 & 5\\
0 & 4 & 2 & 7\\
0 & 0 & 0 & 4\\
\end{block}
\end{blockarray}.
\end{align*}
%
\end{eks}
%
Trappeformen kan reduceres yderligere.
%
\begin{defn}{}{redtrap}
En matrix er på \textbf{reduceret trappeform}, hvis den opfylder følgende kriterier:
\begin{enumerate}[label=(\alph*)]
\item Hvis en søjle har den første ikke-nulindgang i en række, så er alle indgange i søjlen $0$.
\item Den første indgang i hver ikke-nulrække er $1$. 
\end{enumerate}
\end{defn}
\noindent
Den reducerede trappeform af en matrix $A$ noteres  $R$. Et lineært ligningssystem, som ikke står på trappeform, kan laves om til trappeform og derefter reduceret trappeform, hvilket gør løsningen af systemet mere simpel.
Metoden til at gøre dette ses i afsnit \ref{gauss}.
\\
%
\begin{eks}
Betragt matricen $A$.
\begin{align*}
A=
\begin{blockarray}{ccccc}
\begin{block}{[cccc|c]}
1 & 0 & 5 & 0 & 0\\
0 & 1 & 3 & 0 & 0\\
0 & 0 & 0 & 1 & 0\\
0 & 0 & 0 & 0 & 1\\
\end{block}
\end{blockarray}
\end{align*}
%
Jævnfør \ref{defn:redtrap} er $A$ på reduceret trappeform.
\end{eks}
%
%Motiverende tekst
%
%
%%%%%%%%%%%%%%%%%%%%%%%%%%%%%%%%%%%%%%%%%%%%%%%%%%%%%%%%%%%%%
%
\subsection{Pivotindgange}
\begin{defn}{}{}
\textbf{Pivotindgange} betegnes som den første ikke-nulindgang i enhver række i en matrix på trappeform. 
\end{defn}
\noindent
\begin{eks}\label{eks:pivot}
Betragt matrix $A$ fra \ref{eks:trappe}. Pivotindgangene i $A$ er her markeret med blåt.
%
\begin{align*}
A=
\begin{blockarray}{cccc}
\begin{block}{[ccc|c]}
\hlight{2}	& 4			& 8			& 2\\
0			& \hlight{5}& -2		& 5\\
0			& 0			& \hlight{2}& 7\\
0			& 0			& 0			& 0\\
\end{block}
\end{blockarray}.
\end{align*}
%
\end{eks}
Pivotindgange vil fremover blive markeret med blå, som set i \ref{eks:pivot}.
Dette vil blive gjort første gang en matrice bliver lavet til trappeform eller reduceret trappeform.
%\begin{eks}\label{eks:pivot}
%Betragt eksempel \ref{trappe}. Det ses, at ligningssystemet ikke har nogen løsning, da den fjerde ligning ikke har nogen løsning.\\
%\begin{minipage}{0.5\textwidth}
%%
%\begin{align*}
%A=
%\begin{blockarray}{ccccc}
%x_1 & x_2 & x_3 & x_4 & b \\
%\begin{block}{[cccc|c]}
%\hlight{1} & 0 & 5 & 0 & 0\\
%0 & \hlight{1} & 3 & 0 & 0\\
%0 & 0 & 0 & \hlight{1} & 0\\
%0 & 0 & 0 & 0 & \hlight{1}\\
%\end{block}
%\end{blockarray}.
%\end{align*}
%\end{minipage}
%\begin{minipage}{0.5\textwidth}
%\begin{align*}
%x_1+5x_3&=0\\
%x_2+3x_3&=0\\
%x_4&=0\\
%0&\neq 1
%\end{align*}
%\end{minipage}
%%
%\end{eks}
Hvis der er pivotindgang i den sidste søjle, så har ligningssystemet ingen løsning.
%
De variables søjler, der indeholder pivotindgange, kaldes for \textit{faste variable}, mens de, som ikke indholder pivotindgange, kaldes for \textit{frie variable}. 
Ligningssystemet kan løses for de faste variable udtrykt ved de frie variable. 
	\newpage
\subsection{Gauss-elimination}
\label{gauss}
% 1.4
% Sætning 1.4 kan vi bevise den
En algoritme til løsning af lineære ligningssystemer er  \textit{rækkereduktionsalgoritmen}, som reducerer en totalmatrix til reduceret trappeform.
Når den bruges til at løse lineære ligningssystemer, kaldes denne proces for \textit{Gauss-elimination}.
I algoritmen bruges de elementære rækkeoperationer til at omdanne en matrix til trappeform. 
Herefter benyttes rækkeoperationerne til at reducere matricen til reduceret trappeform.
Algoritmen gennemgår således følgende skridt:
%
\begin{enumerate}
\item Find første ikke-nulsøjle fra venstre i matricen.
\item Ved rækkeombytning placeres en ikke-nulindgang øverst i pivotsøjlen.
\item Skab nulindgange under pivotindgangen ved hjælp af rækkeudskiftning.
\item Øverste række markeres som afsluttet og trin $1-4$, hvor afsluttede rækker ignoreres.
Dette gentages, indtil alle rækker er markeret som afsluttet.
\item Alle rækker med pivotindgange skaleres, så alle pivotindgange er lig $1$.
\item Ved rækkeudskiftning sikres nu nulindgange over og under pivotindgangene.
\end{enumerate}
%
%%%%%%%%%%%%%%%%%%%%%%%%%%%%%%%%%%%%%%%%%%%%%%%%%
% Eksempel
\begin{eks}
Betragt ligningssystemet fra \ref{eks:lignsys}
%
\begin{align*}
\begin{array}{rlr}
x_1-x_3-2x_4-8x_5    &=  &-3 \\
-2x_1+x_3+2x_4+9x_5  &=  &5 \\
3x_1-2x_3-3x_4-15x_5 &=  &-9.
\end{array}
\end{align*}
%
Ligningssystemet opskrevet som totalmatrix
%
\begin{equation*}
[A \mid \mathbf{b}] =
\begin{blockarray}{ccccccc}
x_1 & x_2 & x_3 & x_4 & x_5 & b \\
\begin{block}{[ccccc|c]c}
1 & 0 & -1 & -2 & -8 & -3 \\
-2 & 0 & 1 & 2 & 9 & 5 \\
3 & 0 & -2 & -3 & -15 & -9 \\
\end{block}
\end{blockarray}.
\end{equation*}
%
De elementære rækkeoperationer benyttes nu med henblik på at reducere totalmatricen til trappeform.
%
\begin{equation*}
\xrightarrow[R_3 \rightarrow R_3-3R_1]{R_2 \rightarrow R_2+2R_1} 
\begin{blockarray}{ccccccc}
x_1 & x_2 & x_3 & x_4 & x_5 & b \\
\begin{block}{[ccccc|c]c}
  1 & 0 & -1 & -2 & -8 & -3 \\
  0 & 0 & -1 & -2 & -7 & -1 \\
  0 & 0 & 1 & 3 & 9 & 0 \\
\end{block}
\end{blockarray}
\end{equation*}
%
\begin{equation*}
\xrightarrow{R_3 \rightarrow R_3+R_2}
\begin{blockarray}{ccccccc}
x_1 & x_2 & x_3 & x_4 & x_5 & b \\
\begin{block}{[ccccc|c]c}
  \hlight{1} & 0 & -1 & -2 & -8 & -3 \\
  0 & 0 & \hlight{-1} & -2 & -7 & -1 \\
  0 & 0 & 0 & \hlight{1} & 2 & -1 \\
\end{block}
\end{blockarray}
\end{equation*}
%
Algoritmen fortsættes nu med henblik på at reducere totalmatricen fra trappeform til reduceret trappeform.
%
\begin{equation*}
\xrightarrow[R_2 \rightarrow R_2+2R_3]{R_1 \rightarrow R_1+2R_3}
\begin{blockarray}{ccccccc}
x_1 & x_2 & x_3 & x_4 & x_5 & b \\
\begin{block}{[ccccc|c]c}
  1 & 0 & -1 & 0 & -4 & -5 \\
  0 & 0 & -1 & 0 & -3 & -3 \\
  0 & 0 & 0 & 1 & 2 & -1 \\
\end{block}
\end{blockarray}
\end{equation*}
%
\begin{equation*}
\xrightarrow{R_2 \rightarrow -1R_2}
\begin{blockarray}{ccccccc}
x_1 & x_2 & x_3 & x_4 & x_5 & b \\
\begin{block}{[ccccc|c]c}
  1 & 0 & -1 & 0 & -4 & -5 \\
  0 & 0 & 1 & 0 & 3 & 3 \\
  0 & 0 & 0 & 1 & 2 & -1 \\
\end{block}
\end{blockarray}
\end{equation*}
%
\begin{equation*}
\xrightarrow{R_1 \rightarrow R_1+R_2}
\begin{blockarray}{ccccccc}
x_1 & x_2 & x_3 & x_4 & x_5 & b \\
\begin{block}{[ccccc|c]c}
  \hlight{1} & 0 & 0 & 0 & -1 & -2 \\
  0 & 0 & \hlight{1} & 0 & 3 & 3 \\
  0 & 0 & 0 & \hlight{1} & 2 & -1 \\
\end{block}
\end{blockarray}
\end{equation*}
%
Herefter kan følgende løsninger opskrives:
%
\begin{align*}
\begin{array}{rrcll}
x_1-x_5     =&-2   &\iff &x_1   =&x_5-2 \\
x_3 + 3x_5    =&3    &\iff &x_3   =&-3x_5+3 \\
x_4+2x_5    =&-1   &\iff &x_4   =&-2x_5-1.
\end{array}
\end{align*}
%
Løsningen, opskrevet som en parameterfremstilling, er dermed givet ved
%
  \begin{align*}
    \mathbf{x} &= \begin{bmatrix}
           x_{1} \\
           x_{2} \\
           x_{3} \\
           x_{4} \\
           x_{5} \\
         \end{bmatrix} 
         = \begin{bmatrix}
           -2 \\
           0 \\
           3 \\
           -1 \\
           0 \\
         \end{bmatrix}
         +x_2 \begin{bmatrix}
           0 \\
           1 \\
           0 \\
           0 \\
           0 \\
         \end{bmatrix}
         +x_5 \begin{bmatrix}
           1 \\
           0 \\
           -3 \\
           -2 \\
           1 \\
         \end{bmatrix}.
  \end{align*} 
%
\label{eks_gauss}
%
% Frie variable
%
\end{eks}
%
%%%%%%%%%%%%%%%%%%%%%%%%%%%%%%%
%%                           %%
%%    Rang og nullitet       %%
%%                           %%
%%%%%%%%%%%%%%%%%%%%%%%%%%%%%%
%
\subsection{Rang og nullitet}
% 
For at beskrive matricers egenskaber i forhold til pivotindgange er der brug for at definere \textit{rang} og \textit{nullitet}.
%
\begin{defn}{}{rangnull}
\textbf{Rangen} af en matrix $A$ er antallet af pivotindgange, hvilket noteres $\text{rang}(A)$.
\textbf{Nulliteten} af $A$ er antallet af søjler uden pivotindgang, hvilket noteres $\text{null}(A)$.
\end{defn}
%
\noindent
Bemærk, at der for en $m \times n$ matrix gælder, at $\text{rang}(A)+\text{null}(A)=n$.\\
%
\begin{eks}
Med udgangspunkt i matrix $[A \mid \textbf{b}]$ på reduceret trappeform fra \ref{eks_gauss} bliver rang og nullitet i således $\text{rang}(A)=3$ og $\text{null}(A)=2$.
\end{eks}
%
\begin{thm}{}{konsistens}
%
Følgende udsagn er ækvivalente:
%
\begin{enumerate}[label=(\alph*)]
\item Ligningssystemet $A\mathbf{x}=\mathbf{b}$ er konsistent.
\item Vektoren $\mathbf{b}$ er en linearkombination af søjlerne i $A$.
\item Den reducerede trappeform af totalmatricen $[A \mid \mathbf{b}]$ har ingen rækker på formen $[ 0 \ldots 0 \mid d ]$, hvor $d \neq 0$.
\end{enumerate}
%
\end{thm}
%
%
\begin{proof}
%
Lad $A$ være en $m \times n$ matrix og $\mathbf{b} \in \R^m$. Jævnfør \ref{defn:mvp} følger det, at der eksisterer en vektor 
%
\begin{align*}    
       \mathbf{v} = \begin{bmatrix}
		v_{1} \\
        v_{2} \\
        \vdots \\
        v_{n} 
        \end{bmatrix},  
\end{align*}
%
i $\R^n$, der opfylder $A\mathbf{v}=\mathbf{b}$,
hvilket er linearkombinationen
%
$$v_1 \mathbf{a}_1+v_2 \mathbf{a}_2 + \ldots + v_n \mathbf{a}_n = \textbf{b}$$
%
af $A$ og $\textbf{v}$. 
Det følger derfor, at (a) og (b) er ækvivalente.\\\\
%
Herefter bevises, at (a) og (c) er ækvivalente. 
Hvis (c) er falsk, følger det, at
%
$$0 x_1+0 x_2 + \ldots + 0x_n =d,$$
%
hvor $d\neq0$. Da dette ikke har nogen løsning, følger det, at ligningssystemet således bliver inkonsistent. 
Såfremt (c) er sandt, gælder det, at alle rækker i den reducerede trappeform har en løsning, som også er løsning til $A\mathbf{x}=\mathbf{b}$. 
Således er (a) og (c) ækvivalente.
%
\end{proof}
\\
%
\begin{eks}
Med udgangspunkt i \ref{eks_gauss} ses derfor, at ligningssystemet for $A\textbf{x}=\textbf{b}$ er konsistent, da der ikke er pivotindgang i sidste søjle. Såfremt ligningssystemet skulle være inkonsistent, kunne den have formen
%
\begin{align*}
[B \mid \textbf{b} ]=\begin{blockarray}{ccccccc}
x_1 & x_2 & x_3 & x_4 & x_5 & b \\
\begin{block}{[ccccc|c]c}
  \hlight{1} & 0 & 0 & 0 & -1 & -2 \\
  0 & 0 & \hlight{1} & 0 & 3 & 3 \\
  0 & 0 & 0 & 0 & 0 & \hlight{-1} \\
\end{block}
\end{blockarray},
\end{align*}
%
da der eksisterer en pivotindgang i sidste søjle. 
%
\end{eks}
	% Span
	\subsection{Span}
% 1.4

	% Lineær uafhængighed
	\newpage
\noindent
\section{Lineær uafhængighed}
% 1.7
En mængde af vektorer i $\R^n$ er enten \textit{lineært afhængige} eller \textit{lineært uafhængige}.
%
%En mængde af vektorer er \textbf{lineært afhængig}, hvis alle vektorerne kan skrives som en linearkombination af de andre vektorer. 
%Hvis ikke alle vektorerne kan skrives som en linearkombination af de andre vektorer, så er mængden \textbf{lineært uafhængige}. 
%Det kan derfor siges, at hvis mindst en af vektorerne afhænger af de andre vektorer, er samlingen lineært afhængig, og dette leder frem til følgende definition. 
%
%%%%%%%%%%%%%%%%%%%%%%%%%%%%%%%%%%%%%%%%%%%%%%%%%%%%%%%%%%%%%%%%%%%%%%%%%
% 
\begin{defn}{}{linu}
%En mængde af vektorer $\mathbf{u}_1, \mathbf{u}_2, \ldots , \mathbf{u}_k \in \R^n$ kaldes lineært uafhængige, hvis der eksisterer skalarer $x_1, x_2, \ldots , x_k$, sådan at ligningen 
En mængde af vektorerne $\mathbf{u}_1, \mathbf{u}_2, \ldots , \mathbf{u}_k \in \R^n$ er \textbf{lineært uafhængige}, hvis der eksisterer skalarer $x_1, x_2, \ldots , x_k$, sådan at ligningen 
\begin{align*}
x_1\mathbf{u}_1 + x_2\mathbf{u}_2 + \cdots + x_k \mathbf{u}_k = \mathbf{0}
\end{align*}
kun har den trivielle løsning $x_1, x_2, \ldots, x_k = 0$.
Ellers er vektorerne \textbf{lineært afhængige}.
\end{defn}  
%
%%%%%%%%%%%%%%%%%%%%%%%%%%%%%%%%%%%%%%%%%%%%%%%%%
% Eksempel
\begin{eks}\label{fisk}
Givet
\begin{align*}
\S_1 =
\left\{ 
\begin{bmatrix}
           2 \\
           1 \\
\end{bmatrix} 
,
\begin{bmatrix}
           1 \\
           3 \\
\end{bmatrix}
\right\},
\end{align*}
opstilles ligningen  
%
\begin{align*}
x_1 
\begin{bmatrix}
           2 \\
           1 \\
\end{bmatrix} 
+ x_2
\begin{bmatrix}
           1 \\
           3 \\
\end{bmatrix}
= \mathbf{0}.
\end{align*}
%
Denne ligning er kun opfyldt, hvis $x_1=x_2=0$. 
Eftersom der kun eksisterer den trivielle løsning, er vektorerne i $\S_1$ lineært uafhængige.
\\\\
Givet
\begin{align*}
\S_2 =
\left\{ 
\begin{bmatrix}
           2 \\
           1 \\
\end{bmatrix} 
,
\begin{bmatrix}
           4 \\
           2 \\
\end{bmatrix}
\right\},
\end{align*}
%
ville ligningen 
%
\begin{align*}
x_1 
\begin{bmatrix}
           2 \\
           1 \\
\end{bmatrix} 
+ x_2
\begin{bmatrix}
           4 \\
           2 \\
\end{bmatrix}
= \mathbf{0}
\end{align*}
%
%
være opfyldt, hvis $x_1=x_2=0$ eller $x_1=-2$ og $x_2=1$. 
Vektorerne i $\S_2$ er derfor lineært afhængige, fordi der er flere løsninger end den trivielle.
%
\end{eks}
%
%%%%%%%%%%%%%%%%%%%%%%%%%%%%%%%%%%%%%%%%%%%%%%%%%%%%%%%%%%%%%%%%%%%%%%%%%
%
%I eksempel \ref{lineu} viste en af måderne hvorpå man kunne bestemme om et bestemt set $\S$ var lineært afhængigt eller uafhængigt. 
%Der findes flere måder at bestemme dette på. Denne følgende sætning viser en lignende teknik til bestemmelse heraf. 
%
Såfremt en mængde vektorer er lineært uafhængige, eksisterer der en række ækvivalente udsagn, som udtrykt i \ref{thm:mxlinuaf}.
%
%%%%%%%%%%%%%%%%%%%%%%%%%%%%%%%%%%%%%%%%%%%%%%%%%%%%%%%%%%%%%%%%%%%%%%%%%
%
\begin{thm}{}{mxlinuaf}
%
Lad $A$ være en $m \times n$ matrix.
Følgende udsagn er ækvivalente:
%
\begin{enumerate}[label=(\alph*)]
\item $A$'s søjler er lineært uafhængige. 
\item Ligningssystemet $A\mathbf{x}=\mathbf{b}$ har højst én løsning for alle $\mathbf{b} \in \R^m$.
\item $\text{null}(A)=0$.
\item $\text{rang}(A)=n$.
\item Søjlerne i $A$ på reduceret trappeform er standardvektorer i $\R^m$.
\item $A\mathbf{x}=\mathbf{0}$ har kun løsningen $\mathbf{x}=\mathbf{0}$.                                                                                                                                                                                                                                                                                                                         
\item Der er pivotindgange i hver søjle i $A$. 
\end{enumerate}
%
\end{thm}
%
%%%%%%%%%%%%%%%%%%%%%%%%%%%%%%%%%%%%%%%%%%%%%%%%%%%%%%%%%%%%%%%%%%%%%%%%%
%
\begin{proof}
%Forudsagt? Mangler der en ref til et tidligere bevis her, eller er I spåkællinger? (nedenstående) 
Lad $A$ være en $m \times n$ matrix.
Det er følger direkte af \ref{defn:linu}, at (a) og (f), samt (f) og (g) er ækvivalente, så for at færdiggøre beviset for denne sætning skal det vises, at (b) $\rightarrow$ (c), (c) $\rightarrow$ (d), (d) $\rightarrow$ (e), (e) $\rightarrow$ (f) og (f) $\rightarrow$ (b).
\\\\
%Er de pile korrekt notation, eller skal vi bruge \equiv? (ovenstående)
% (b) $\rightarrow$ (c). 
Idet $\mathbf{0}$ er en løsning til $A\mathbf{x}=\mathbf{b}$, så forudsætter (b), at $A\mathbf{x}=\mathbf{b}$ ikke har ikke-nul løsninger, da den generelle løsning ingen frie variable har. 
Eftersom antallet af frie variable er $0$, og $\text{null}(A)$ dermed er nul, er (b) $\rightarrow$ (c) hermed bevist. 
\\\\
%Der skal ikke komma efter "idet". 
% (c) $\rightarrow$ (d).
Idet $\text{rang}(A)+\text{null}(A)=n$, så er (c) $\rightarrow$ (d) bevist.
\\\\
%
% (d) $\rightarrow$ (e).
Hvis $\text{rang}(A)=n$, så indeholder enhver søjle i $A$ en pivotindgang, hvormed den reducerede trappeform udelukkende indeholder standardvektorer. 
Disse standardvektorer er nødvendigvis forskellige, da hver søjle indeholder den første ikke-nul indgang i hver række. 
Dette beviser (d) $\rightarrow$ (e).
%
\\\\
% (e) $\rightarrow$ (f).
Lad $A_R$ være den reducerede trappeform af matricen $A$. Hvis søjlerne i $A_R$ er forskellige standardvektorer i $\R^m$, så er $A_R= [ \mathbf{e}_1, \mathbf{e}_2, \ldots, \mathbf{e}_n]$. 
Det er givet, at den eneste løsning til $A_R\mathbf{x}=\mathbf{0}$ er $\mathbf{0}$, og idet $A\mathbf{x}=\mathbf{0}$ er ækvivalent med $A_R\mathbf{x}=\mathbf{0}$, så er $\mathbf{0}$ eneste løsning til $A\mathbf{x}=\mathbf{0}$. 
Dette beviser (e) $\rightarrow$ (f).
\\\\
%
% (f) $\rightarrow$ (b).
Lad $\mathbf{u}$ og $\mathbf{v}$ være løsninger til $A\mathbf{x}=\mathbf{b}$.
% Sætning 1.3
Jævnfør \ref{thm:mxvpro} følger det, at $A(\mathbf{u}-\mathbf{v})=A\mathbf{u}-A\mathbf{v}=\mathbf{b}-\mathbf{b}=\mathbf{0}$, hvilket tilsvarer $\mathbf{u}-\mathbf{v}=\mathbf{0}
\rightarrow  
\mathbf{u} =\mathbf{v}$. Dette beviser (f) $\rightarrow$ (b).
%\\\\
%
%Lad $\mathbf{b}=\mathbf{0}$. 
%Ligningssystemet $A\mathbf{x}=\mathbf{b}$ har altid løsningen $\mathbf{x}=\mathbf{0}$, der derfor er den eneste løsning til $A\mathbf{x}=\mathbf{0}$, hvilket beviser (b) $\rightarrow$ (e).
%
\end{proof}
\\
%
%%%%%%%%%%%%%%%%%%%%%%%%%%%%%%%%%%%%%%%%%%%%%%%%%%%%%%%%%%%%%%%%%%%%%%%%%
%
Ligeså findes der en egenskab for lineært afhængige vektorer, som er belyst i \ref{thm:inspektion}.
%
%%%%%%%%%%%%%%%%%%%%%%%%%%%%%%%%%%%%%%%%%%%%%%%%%%%%%%%%%%%%%%%%%%%%%%%%%
%
\begin{thm}{}{inspektion}
%
Vektorerne $\mathbf{u}_1,\mathbf{u}_2, \ldots ,\mathbf{u}_k \in\R^n$ er lineært afhængige, 
hvis og kun hvis $\mathbf{u}_1=\mathbf{0}$, eller hvis der findes $i \geq 2$, således $\mathbf{u}_i \in \text{span} \{ \mathbf{u}_1,\mathbf{u}_2, \ldots ,\mathbf{u}_{i-1} \}$.
%
\end{thm}
%
%%%%%%%%%%%%%%%%%%%%%%%%%%%%%%%%%%%%%%%%%%%%%%%%%%%%%%%%%%%%%%%%%%%%%%%%%
%
\begin{proof}
% 
Antag, at $\mathbf{u}_1 \neq \mathbf{0}$, da løsningen ellers er triviel. 
Der findes skalarer, som ikke alle er $0$, således at
% 
\begin{align*}
x_1 \mathbf{u}_1 + x_2 \mathbf{u}_2 + \ldots + x_k \mathbf{u}_k = \mathbf{0}.
\end{align*}
%
Lad $i$ være det største index $i$, hvor $x_i \neq 0$, sådan at 
\begin{align*}
x_1 \mathbf{u}_1 + x_2 \mathbf{u}_2 + \ldots + x_i \mathbf{u}_i = \mathbf{0}.
\end{align*}
%Synes det ovenstående lyder lidt mærkeligt (det med i, i og x_i). 
Bemærk, at $i\geq 2$, da udtrykket ellers ville blive reduceret til
%
\begin{align*}
x_1\textbf{u}_1=\mathbf{0},
\end{align*}
%
hvilket går imod tidligere antagelser.
Løses udtrykket i forhold til $\mathbf{u}_i$, ses det, at
\begin{align*}
\mathbf{u}_i = - \frac{x_1}{x_i} \mathbf{u}_1 - \frac{x_2}{x_i} \mathbf{u}_2 - \cdots - \frac{x_{i-1}}{x_i} \mathbf{u}_{i-1}.
\end{align*}
Hvilket vil sige, at 
\begin{align*}
\mathbf{u}_i \in \text{span}\{ \mathbf{u}_1,\mathbf{u}_2, \ldots ,\mathbf{u}_{i-1} \}.
\end{align*}
%
Antag nu, at $\mathbf{u}_i \in \text{span}\{ \mathbf{u}_1,\mathbf{u}_2, \ldots ,\mathbf{u}_{i-1} \}$ for $i \geq 2$.
Da $\text{span}\{ \mathbf{u}_1,\mathbf{u}_2, \ldots ,\mathbf{u}_{i-1} \}$ er mængden af alle linearkombinationer af vektorerne $\mathbf{u}_1,\mathbf{u}_2, \ldots ,\mathbf{u}_{i-1}$, eksisterer der skalarer $x_1, x_2, \ldots x_{i-1}$, således at
%
\begin{align*}
\textbf{u}_i=x_1\mathbf{u}_1 + x_2 \mathbf{u}_2 + \cdots + x_{i-1}\mathbf{u}_{i-1}.
\end{align*}
%
For $k $ haves, at $\mathbf{u}_1,\mathbf{u}_2, \ldots ,\mathbf{u}_k$ er lineært afhængige, da 
\begin{align*}
\textbf{u}_1 -x_1\textbf{u}_1 - x_2\textbf{u}_2 - \cdots - x_{i-1}\textbf{u}_{i-1} = \textbf{0}
\end{align*} 
er en løsning anderledes fra den trivielle løsning.
%
\end{proof}
\\
%
%%%%%%%%%%%%%%%%%%%%%%%%%%%%%%%%%%%%%%%%%%%%%%%%%%%%%%%%%%%%%%%%%%%%%%%%%
%
Det kan i nogle tilfælde afgøres ved inspektion, hvorvidt en mængde vektorer $\mathbf{u}_1,\mathbf{u}_2,\ldots ,\mathbf{u}_k$ er lineært afhængige. 
En mængden af vektorer er lineært afhængige, hvis:
%
\begin{itemize}
\item $\mathbf{0}$ er i mængden.
\item To af vektorerne er parallelle.
\item Der er flere vektorer end komponenter i vektorerne. 
\end{itemize}
%
Med udgangspunkt i \ref{fisk} ses det, at de to vektorer i $\S_2$ er parallelle og derfor er lineært afhængige. 
\\\\
Hvis det ikke ved inspektion kan afgøres, om en matrix er lineært afhængig eller lineært uafhængig, kan det undersøges ved at opstille en totalmatricen bestående af mængden af vektorer, hvor matricen derefter via rækkeopperationer omdannes til trappeform. 
Hvis der er pivotindgange i alle søjlerne, er mængden af vektorer lineært uafhængige. Hvis der er frie variable, er mængden af vektorer derimod lineært afhængige.
\\\\
\newpage
\noindent
%
% Eksempel
% Virker stadig ufærdigt, vi snakker om det. 
\begin{eks}\label{lineu}
Lad
\begin{align*}
\S_1 &= \left\{
\begin{bmatrix}
           4 \\
           1 \\
           1 \\
\end{bmatrix}
,
\begin{bmatrix}
           -3 \\
           0 \\
           1 \\
\end{bmatrix}
,
\begin{bmatrix}
           1 \\
           -2 \\
           1 \\
\end{bmatrix}
\right\}.
\end{align*}
\noindent
%
Hvis vektorerne er lineært uafhængige, er den trivielle løsning den eneste løsning til
%
\begin{align*}
x_1
\begin{bmatrix}
           4 \\
           1 \\
           1 \\
\end{bmatrix}
+ x_2
\begin{bmatrix}
           -3 \\
           0 \\
           1 \\
\end{bmatrix}
+ x_3
\begin{bmatrix}
           1 \\
           -2 \\
           1 \\
\end{bmatrix}
=\mathbf{0}.
\end{align*}
%
Ved hjælp af rækkeoperationerne omskrives totalmatricen til reduceret trappeform:
%
\noindent
\begin{align*}
[A \mid \textbf{b}]=
\begin{blockarray}{cccc}
x_1 & x_2 & x_3 & b \\
\begin{block}{[ccc|c]}
4 & -3 & 1 & 0\\
1 & 0 & -2 & 0 \\
1 & 1 & 1 & 0 \\
\end{block}
\end{blockarray}
\xrightarrow[R_2 \rightarrow R_2-1R_1]{R_1 \rightarrow \frac{1}{4}R_1} 
\begin{blockarray}{cccc}
x_1 & x_2 & x_3 & b \\ 
\begin{block}{[ccc|c]}
1 & -\frac{3}{4} & \frac{1}{4} & 0\\
0 & \frac{3}{4} & -\frac{9}{4} & 0 \\
1 & 1 & 1 & 0 \\
\end{block}
\end{blockarray} \\
\xrightarrow[R_2 \rightarrow \frac{4}{3} R_2]{R_3 \rightarrow R_3-1R_1} 
\begin{blockarray}{cccc}
x_1 & x_2 & x_3 & b \\
\begin{block}{[ccc|c]}
1 & -\frac{3}{4} & \frac{1}{4} & 0\\
0 & 1 & -3 & 0\\
0 & \frac{7}{4} & \frac{3}{4} & 0 \\
\end{block}
\end{blockarray}
\xrightarrow[R_3 \rightarrow \frac{1}{6} R_3]{R_3 \rightarrow R_3-\frac{7}{4}R_2} 
\begin{blockarray}{cccc}
x_1 & x_2 & x_3 & b \\
\begin{block}{[ccc|c]}
\hlight{1} & \frac{-3}{4} & \frac{1}{4} & 0\\
0 & \hlight{1} & -3 & 0\\
0 & 0 & \hlight{1} & 0 \\
\end{block}
\end{blockarray} \\
\xrightarrow[R_1 \rightarrow R_1-\frac{1}{4}R_3]{R_2 \rightarrow R_2+3R_3} 
\begin{blockarray}{cccc}
x_1 & x_2 & x_3 & b \\
\begin{block}{[ccc|c]}
1 & \frac{-3}{4} & 0 & 0\\
0 & 1 & 0 & 0\\
0 & 0 & 1 & 0 \\
\end{block}
\end{blockarray} 
\xrightarrow[]{R_1 \rightarrow R_1+\frac{3}{4}R_2} 
\begin{blockarray}{cccc}
x_1 & x_2 & x_3 & b \\
\begin{block}{[ccc|c]}
\hlight{1} & 0 & 0 & 0\\
0 & \hlight{1} & 0 & 0\\
0 & 0 & \hlight{1} & 0 \\
\end{block}
\end{blockarray} \\
\end{align*}
%
Eftersom totalmatricen ikke har nogle frie variable, har den kun en løsning, og er derfor lineært uafhængig.
\\\\
%
\noindent
Lad nu
%
\begin{align*}
\S_2 &= \left\{
\begin{bmatrix}
           4 \\
           1 \\
           -2 \\
\end{bmatrix}
,
\begin{bmatrix}
           -3 \\
           0 \\
           1 \\
\end{bmatrix}
,
\begin{bmatrix}
           1 \\
           -2 \\
           1 \\
\end{bmatrix}
\right\}.
\end{align*}
%
Nu afgøres, om den trivielle løsning er den eneste løsning til
%
\begin{align*}
x_1+
\begin{bmatrix}
           4 \\
           1 \\
           -2 \\
\end{bmatrix}
+ x_2
\begin{bmatrix}
          -3 \\
           0 \\
           1 \\
\end{bmatrix}
+ x_3
\begin{bmatrix}
           1 \\
           -2 \\
           1 \\
\end{bmatrix}
=\textbf{0}.
\end{align*}
%
Ved hjælp af rækkeoperationer omskrives totalmatricen til trappeform:
%
\begin{align*}
[B \mid \textbf{b}]=&
\begin{blockarray}{cccc}
x_1 & x_2 & x_3 & b \\
\begin{block}{[ccc|c]}
4 & -3 & 1 & 0\\
1 & 0 & -2 & 0 \\
-2 & 1 & 1 & 0 \\
\end{block}
\end{blockarray}\\
\xrightarrow[R_2\rightarrow R_2-4R_1]{R_1 \leftrightarrow R_2}&
\begin{blockarray}{cccc}
x_1 & x_2 & x_3 & b \\
\begin{block}{[ccc|c]}
1 & 0 & -2 & 0\\
0 & -3 & 9 & 0\\
-2 & 1 & 1 & 0 \\
\end{block}
\end{blockarray}\\
\xrightarrow[R_3 \rightarrow R_3+\frac{1}{3}R_2]{R_3 \rightarrow R_3+2R_1}&
\begin{blockarray}{cccc}
x_1 & x_2 & x_3 & b \\
\begin{block}{[ccc|c]}
\hlight{1} & 0 & -2 & 0\\
0 & \hlight{-3} & 9 & 0 \\
0 & 0 & 0 & 0 \\
\end{block}
\end{blockarray}
\end{align*}
%
Dette betyder, at der findes en ikke-triviel linearkombination, som resulterer i nulvektoren.
Derfor er der mindst én ikke-triviel linearkombination, som opfylder $A\textbf{x}=\textbf{0}$, og derfor er vektormængden lineært afhængig.
\end{eks}
%Der bliver stadig skrevet "sæt" rigtig meget.. :-( 
%%%%%%%%%%%%%%%%%%%%%%%%%%%%%%%%%%%%%%%%%%%%%%%%%%%%%%%%%%%%%%%%%%%%%%%%%
%
% SKAL MÅSKE FLYTTES
I forbindelse med en given mængde vektorer kan færre vektorer udvælges for at få en mere effektiv beskrivelse af spannet. Dette gøres ved at finde spannet af pivotsøjlerne i koefficientmatricen.
\\\\
%
%%%%%%%%%%%%%%%%%%%%%%%%%%%%%%%%%%%%%%%%%%%%%%%%%
% Eksempel
\begin{eks}
Givet 
\begin{align*}
\S &= \left\{
\begin{bmatrix}
           1 \\
           2 \\
           0 \\
\end{bmatrix}
,
\begin{bmatrix}
           -1 \\
           2 \\
           4 \\
\end{bmatrix}
,
\begin{bmatrix}
           3 \\
           0 \\
           1 \\
\end{bmatrix}
,
\begin{bmatrix}
           1 \\
           0 \\
           5 \\
\end{bmatrix}
\right\},
\end{align*}
% 
%
opstilles koefficientmatricen, som omdannes til reduceret trappeform. 
%
\begin{align*}
&\begin{blockarray}{ccccc}
\begin{block}{[cccc]c}
  1 & -1 & 3 & 1 \\
  2 & 2 & 0 & 0 \\
  0 & 4 & 1 & 5 \\
\end{block}
\end{blockarray}
%
& \xrightarrow{R_2 \rightarrow R_2-2R_1}&
\begin{blockarray}{ccccc}
\begin{block}{[cccc]c}
  1 & -1 & 3 & 1 \\
  0 & 4 & -6 & -2 \\
  0 & 4 & 1 & 5 \\
\end{block}
\end{blockarray} \\
%
\xrightarrow{R_3 \rightarrow R_3-R_2}&
\begin{blockarray}{ccccc}
\begin{block}{[cccc]c}
  \hlight{1} & -1 & 3 & 1 \\
  0 & \hlight{4} & -6 & -2 \\
  0 & 0 & \hlight{7} & 7 \\
\end{block}
\end{blockarray} 
%
&\xrightarrow{R_3 \rightarrow \frac{1}{7} R_3}&
\begin{blockarray}{ccccc}
\begin{block}{[cccc]c}
  1 & -1 & 3 & 1 \\
  0 & 4 & -6 & -2 \\
  0 & 0 & 1 & 1 \\
\end{block}
\end{blockarray} \\
%
\end{align*}
\begin{align*}
%
\xrightarrow{R_2 \rightarrow R_2+6R_3}&
\begin{blockarray}{ccccc}
\begin{block}{[cccc]c}
  1 & -1 & 3 & 1 \\
  0 & 4 & 0 & 4 \\
  0 & 0 & 1 & 1 \\
\end{block}
\end{blockarray} 
%
&\xrightarrow{R_1 \rightarrow R_1-3R_3}&
\begin{blockarray}{ccccc}
\begin{block}{[cccc]c}
  1 & -1 & 0 & -2 \\
  0 & 4 & 0 & 4 \\
  0 & 0 & 1 & 1 \\
\end{block}
\end{blockarray} \\
%
\xrightarrow{R_2 \rightarrow \frac{1}{4} R_2}&
\begin{blockarray}{ccccc}
\begin{block}{[cccc]c}
  1 & -1 & 0 & -2 \\
  0 & 1 & 0 & 1 \\
  0 & 0 & 1 & 1 \\
\end{block}
\end{blockarray} 
%
&\xrightarrow{R_1 \rightarrow R_1+R_2}&
\begin{blockarray}{ccccc}
\begin{block}{[cccc]c}
  \hlight{1} & 0 & 0 & -1 \\
  0 & \hlight{1} & 0 & 1 \\
  0 & 0 & \hlight{1} & 1 \\
\end{block}
\end{blockarray}.
%
\end{align*}
%
Hermed er den sidste vektor i spannet en linearkombination af de andre vektorer:
%
  \begin{align*}
         -1 \begin{bmatrix}
           1 \\
           2 \\
           0 \\
         \end{bmatrix}
         +1
         \begin{bmatrix}
           -1 \\
           2 \\
           4 \\
         \end{bmatrix}
          +1
         \begin{bmatrix}
           3 \\
           0 \\
           1 \\
         \end{bmatrix}
         =
         \begin{bmatrix}
           1 \\
           0 \\
           5 \\
         \end{bmatrix}.
  \end{align*} 
%
Dermed haves, at
%
\begin{align*}
\text{span}\{\S \} &=\text{span}
\left\{
\begin{bmatrix}
           1 \\
           2 \\
           0 \\
\end{bmatrix}
,
\begin{bmatrix}
           -1 \\
           2 \\
           4 \\
\end{bmatrix}
,
\begin{bmatrix}
           3 \\
           0 \\
           1 \\
\end{bmatrix}
\right\}.
\end{align*}
\end{eks}
%
%%%%%%%%%%%%%%%%%%%%%%%%%%%%%%%%%%%%%%%%%%%%%%%%%%%%%%%%%%%%%%%%%%%%%%%%%
%
%\begin{eks}\label{lineu}
%Lad et sæt af vektorer være
%\begin{align*}
%\S &= \left\{
%\begin{bmatrix}
%           2 \\
%           1 \\
%\end{bmatrix}
%,
%\begin{bmatrix}
%           3 \\
%           2 \\
%\end{bmatrix}
%\right\}.
%\end{align*}
%\noindent
%Bestem om vektorerne er lineært uafhængige ved ligningen, altså hvis $x_1$ og $x_2 = 0$
%\begin{align*}
%x_1+
%\begin{bmatrix}
%           2 \\
%           1 \\
%\end{bmatrix}
%+ x_2
%\begin{bmatrix}
%           3 \\
%           2 \\
%\end{bmatrix}
%=0 
%\end{align*}
%%
%Ligningerne udskrives
%\noindent
%\begin{align*}
%2x_1+3x_2=0\\
%x_1+2x_2=0\\
%\\
%x_1+3/2x_2=0\\
%x_1+2x_2=0\\
%\\
%x_1+2x_2=0 - x_1+3/2x_2=0\\
%1/2c_2=0\\
%c_2=0\\
%c_1=0
%\end{align*}
%%
%Da $c_1$ og $c_2 = 0$, så er sættet $\S$ lineært uafhængigt.
%\\
%\\
%\noindent
%Lad et sæt af vektorer være 
%%
%\begin{align*}
%\S &= \left\{
%\begin{bmatrix}
%           2 \\
%           1 \\
%\end{bmatrix}
%,
%\begin{bmatrix}
%           3 \\
%           2 \\
%\end{bmatrix}
%,
%\begin{bmatrix}
%           1 \\
%           2 \\
%\end{bmatrix}
%\right\}.
%\end{align*}
%\noindent
%Bestem om vektorerne er lineært afhængige ved ligningen, altså hvis $x_1$, $x_2$ og $x_3 \neq 0$
%\begin{align*}
%x_1+
%\begin{bmatrix}
%           2 \\
%           1 \\
%\end{bmatrix}
%+ x_2
%\begin{bmatrix}
%           3 \\
%           2 \\
%\end{bmatrix}
%+ x_3
%\begin{bmatrix}
%           1 \\
%           2 \\
%\end{bmatrix}
%=0.
%\end{align*}
%%
%Ligningerne udskrives 
%%
%\begin{align*}
%2x_1+3x_2+c_3=0\\
%x_1+2x_2+2c_3=0 
%\end{align*}
%\noindent
%Nu findes en linearkombination, som ikke er nul. 
%Vi sætter tilfældigt $x_3 = -1$ og forsøger om ligningen kan løses
%%
%\begin{align*}
%2x_1+3x_2-1=0\\
%x_1+2x_2-2=0\\
%\\
%2x_1+3x_2-1=0\\
%2x_1+4x_2-4=0\\
%\\	 
%2x_1+3x_2-1=0 - 2x_1+4x_2-4=0 \\
%-x_2+3=0\\
%-x_2=-3\\
%x_2=3\\
%\\
%x_1+6+2\\
%x_1+4=0\\
%x_1=-4
%\end{align*}
%%
%Linearkombinationen er derfor
%%
%\begin{align*}
%-4+
%\begin{bmatrix}
%           2 \\
%           1 \\
%\end{bmatrix}
%+ 3
%\begin{bmatrix}
%           3 \\
%           2 \\
%\end{bmatrix}
%+ -1
%\begin{bmatrix}
%           1 \\
%           2 \\
%\end{bmatrix}
%=0
%\end{align*}
%%
%Derfor har vi minimum fundet én linearkombination som er ikke-nul, som giver os nulvektoren, og derfor er det bevist at vektorsættet lineært afhængig.
%\end{eks}
%


	\input{incl/main/linear_algebra/lineaer_afbilding}
	% Inversematricer
	\section{Elementære og inverse matricer}
At multiplicere en matrix $A$ med en $n \times n$ \textit{elementær matrix} $E$ svarer til at udføre én rækkeoperation på $A$. 
Dermed kan enhver rækkeoperation udføres gennem matrixmultiplikation. 
En $n \times n$ matrix $E$ er elementær, hvis én rækkeoperation udført på $I_n$ resulterer i $E$. 
%Eksempel? 
\\\\
Lad nu $A$ være en $m \times n$ matrix og $E$ en elementær $m \times m$ matrix, der er resultatet af én rækkeoperation udført på $I_m$. 
Matrixproduktet $EA$ kan opnås ved at udføre én rækkeoperation på $A$. 
Enhver elementær matrix er desuden \textit{invertibel}, og dens \textit{invers} er også en elementær matrix. 
%
\begin{defn}{}{mxinvers}
En $n \times n$ matrix $A$ er \textbf{invertibel}, hvis en $n \times n$ matrix $B$ eksisterer, så $AB=BA=I_n$. 
\end{defn}
\noindent
%
Det er endvidere relevant, hvorvidt inversen er entydig.
\begin{thm}{}{mxinversent}
Såfremt en matrix $A$ er invertibel findes der en entydig matrix $B$, som er dennes invers.
\end{thm}
%
\begin{proof}
Antag at der findes to matricer  $B$ og $C$, som begge har størelsen $n \times n$, hvorom det gælder, at disse er inverse til $n \times n$ matricen $A$.
Det gælder derfor fra definition \ref{defn:mxinvers}, at $A$ er inversen til $B$ og $C$.
Derfor haves at $AB=BA=I_{n}$, samt at $AC=CA=I_{n}$. 
Fra sætning \ref{thm:mxmulti} haves derfor at:
\begin{align*}
B=BI_n&=B(AC) \\
&=(BA)C \\
&=I_nC=C
\end{align*} 
Altså er $B=C$ og $A$'s invers er derfor entydig, denne kaldes $A^{-1}$.
\end{proof} \\
%
Såfremt $A$ er en invertibel $n \times n$ matrix, gælder det for ethvert $\mathbf{b} \in \R^n$, at $A\textbf{x}=\mathbf{b}$ har den entydige løsning $A^{-1}\mathbf{b}$.
Inverse matricer har en række egenskaber, hvoraf nogle centrale er formuleret i sætning \ref{thm:mxinvers}. 
%
\begin{thm}{}{mxinvers}
Givet to $m \times n$ matricer $A$ og $B$, så gælder, at
\begin{enumerate}[label=(\alph*)]
\item Hvis $A$ er invertibel, så er $A^{-1}$ invertibel, og $(A^{-1})^{-1}=A$.
\item Hvis $A$ og $B$ er invertible, så er $AB$ invertibel, og $(AB)^{-1}=A^{-1}B^{-1}$.
\item Hvis $A$ er invertibel, så er $A^T$ invertibel, og $(A^T)^{-1}=(A^{-1})^T$.
\end{enumerate}
\end{thm}
%
%
\begin{proof}
(a)
Udsagnet følger af definition \ref{defn:mxinvers}. 
\\\\
(b) Antag, at $A$ og $B$ er invertibel matricer, så $(AB)(B^{-1}A^{-1})=A(BB^{-1})A^{-1}=AI_nA^{-1}=AA^{-1}=I_n$. 
Det kan gennem lignende operationer vises at $(B^{-1}A^{-1})(AB)=I_n$.
Dermed er $AB$ en invertibel matrix, hvis invers er $A^{-1}B^{-1}$.
Dermed er $AB^{-1}=A^{-1}B^{-1}$.
\\\\
(c) Antag at A er invertibel. 
Det følger fra definition \ref{defn:mxinvers}, at $A^{-1}A=I_n$. 
Jævnfør sætning \ref{thm:mxmulti}(g) følger det hermed, at $A^T(A^{-1})^T=(A^{-1}A)^T=I_n^T=I_n$. 
Ligeledes kan det vises, at $(A^{-1})^TA^T=I_n$. 
Derfor er $A^T$ en invertibel matrix med inversen $(A{-1})^T$ og $(A^T)^{-1}A=(A^{-1})^T$.
\end{proof}
%
%NOGET OM ELEMENTÆRE MATRICER OG INVERTIBILITET
%
\begin{thm}{}{eleinv}
Lad $A$ være en $m \times n$ matrix, og $R$ den reducerede trappeform af $A$. 
Så eksisterer der en invertibel $m \times m$ matrix $P$, så $PA=R$.
\end{thm}
%
\begin{proof}
\textbf{Beviset kommer senere }
\end{proof}
%
\\
Sætning \ref{thm:eleinv} medfører desuden korollar \ref{kor:eleinv}.
%måske noget mere her (at det er vigtigt for rækkeoperationer?).  
\begin{kor}{}{eleinv}
Matrixligningerne $A\textbf{x}=\textbf{b}$ og $R\textbf{x}=\textbf{c}$ har samme løsninger, hvor 
$[R \text{   }| \textbf{c}]$
er den reducerede trappeform af den udvidede matrix 
$[A \text{   }| \textbf{b}].
$
\end{kor}\noindent
%
%\begin{proof}
%Lad $P$ være en invertibel matrix, så 
%$P[A \text{   }| \textbf{b}]
%=
%[R \text{   }| \textbf{c}]
%$.
%Så gælder det, at 
%$$
%P
%\begin{bmatrix}
%A & \textbf{b}
%\end{bmatrix}
%=
%\begin{bmatrix}
%PA & P\textbf{b}
%\end{bmatrix}
%=
%\begin{bmatrix}
%R & \textbf{c}
%\end{bmatrix}
%\text{, }
%$$
%hvormed $PA=R$ og $P\textbf{b}=\textbf{c}$.
%Heraf følger, at $A=P^{-1}R$ og $\textbf{b}=P^{-1}\textbf{c}$, da $P$ er invertibel. 
%\\\\
%Antag, at $\textbf{v}$ er en løsning til $A\textbf{x}=\textbf{b}$. Så er
%$$
%R\textbf{v}=(PA)\textbf{v}=P(A\textbf{v})=P\textbf{b}=\textbf{c},
%$$
%hvormed $\textbf{v}$ er en løsning til $R\textbf{x}=\textbf{c}$. 
%Antag nu i stedet, at $\textbf{v}$ er en løsning til $R\textbf{x}=\textbf{c}$. Så er
%$$
%A\textbf{v}=(P^{-1}R)\textbf{v}=P^{-1}(R\textbf{v})=P^{-1}\textbf{c}=\textbf{b},
%$$
%hvormed $\textbf{v}$ er en løsning til $A\textbf{x}=b$. 
%Dermed er det bevist, at $A\textbf{x}=b$ og $R\textbf{x}=c$ har samme løsning.
%%Beviset for korollar \ref{kor:eleinv} følger af beviset for sætning \ref{thm:eleinv} med tilføjelsen, at hvis $\textbf{b}=\textbf{0}$, så er $\textbf{c}=\textbf{0}$, hvilket resulterer i nulvektoren i begge udtryk.
%\end{proof}
%\\
%
%
%HER SAD DEF 3.8 FØR
%%%%%%%%%%%%%%%%%%%%%%%%%%%%%
Det er dog ikke altid, der findes en matrix $B$, så kriterierne i definition \ref{defn:mxinvers} opfyldes. 
%
\begin{thm}{}{mxinvertibel}
Lad $A$ være en $n \times n$ matrix. 
$A$ er invertibel, hvis og kun hvis den reducerede trappeform af $A$ er $I_n$.
\end{thm}
%
\begin{proof}
Antag, at $n\times n$ matricen $A$ er invertibel, og at der findes en vektor \textbf{u} i $\R^n$, så $A\textbf{v}=\textbf{0}$. 
Da $A$ er en invertibel matrix, gælder det for ethvert $\mathbf{b} \in \R^n$, at $A\textbf{x}=\mathbf{b}$ har løsningen $A^{-1}\mathbf{b}$ og derfor er $\textbf{v}=A^{-1}\textbf{0}=\textbf{0}$. 
Den eneste løsning til problemet er derfor $\textbf{0}$. 
Såfremt $A$ er invertibel, følger det derfor fra sætning \ref{thm:mxlinuaf}, at $\text{ rang}(A)=n$.
Det følger derfor, at $A$ på reduceret trappeform er $I_n$.
\\\\
Antag nu, at den reducerede trappeform af $A$=$I_n$. 
Jævnfør sætning \ref{thm:mxinvertibel} eksister der en invertibel $n \times n$ matrix $P$, hvorom det gælder, at $PA=I_n$. 
%
\begin{align*}
A=I_nA=(P^{-1}P)A=P^{-1}(PA)=P^{-1}.
\end{align*}
%
Fra sætning \ref{thm:mxinvers} fremgår det, at $P^{-1}$ er en invertibel matrix, og det er hermed bevist, at $A$ er invertibel, da $A=P^{-1}$.
\end{proof}
%%%%%%%%%%%%%%%%%%%%%%%%%%%%%%
%find begin example
\\
Det er fordelagtigt at benytte algoritmer til at invertere matricer.
Jævnfør definition \ref{defn:mxinvers} opfylder en $n \times n$ matrices inverse matrix $AA^{-1}=I_n$.
Med dette kan ligningen $A[\textbf{x}_1\text{  }\textbf{x}_2\text{  }\ldots\text{  }\textbf{x}_n ] = [\textbf{e}_1\text{  }\textbf{e}_2\text{  }\ldots\text{  }\textbf{e}_n ]$ opstilles.
Da den inverse matrix er entydig kan søjlerne findes ved at løse $A\textbf{x}_i = \textbf{e}_i$.
Det gøres ved Gauss-elimination på $[A\textbf{ }|\textbf{ }\textbf{e}_i]$.
Kombinationen af rækkeoperationer for at få $A$ på reduceret trappeform er kun afhængig af $A$.
Dette betyder, at hver ligning løses med samme rækkeoperationer og alle søjler kan findes samtidigt.
Dette gøres med Gauss-elimination på $[A\textbf{ }|\textbf{ }I_n]$, hvor resultatet bliver $[I_n\textbf{ }|\textbf{ }A^{-1}]$, hvis $A$ er invertibel.\\
%
%Lad $A$ være en $n \times n$ matrix. 
%Givet et input $[A \text{   } I_n]$ skal algoritmen udføre elementære rækkeoperationer på $A$, således algoritmens output bliver $[R \text{   } B]$, hvor $R$ er $A$ på reduceret trappeform. 
%Hvis $R=I_n$, er $A$ invertibel og $B=A^{-1}$.
%Hvis $R \neq I_n$, er $A$ ikke invertibel.\\
%
\begin{eks}
\label{eks:fisk5}
Lad
\begin{align*}
A=\begin{blockarray}{ccc}
\begin{block}{[ccc]}
1 & 1 & 1  \\
1 & 2 & 2  \\
1 & 3 & 2  \\
\end{block}
\end{blockarray}.
\end{align*}
%
For at finde dennes invers $A^{-1}$ opstilles totalmatricen 
\begin{align*}
[A \text{   } | \text{ } I_3] =\begin{blockarray}{cccccc}
\begin{block}{[ccc|ccc]}
1 & 1 & 1 & 1 & 0 & 0 \\
1 & 2 & 2 & 0 & 1 & 0 \\
1 & 3 & 2 & 0 & 0 & 1 \\
\end{block}
\end{blockarray}.
\end{align*}
Ved at udføre rækkeoperationer omdannes $A$ til reduceret trappeform.
\begin{align*}
[A \text{   } | \text{ } I_3] = &
\begin{blockarray}{cccccc}
\begin{block}{[ccc|ccc]}
1 & 1 & 1 & 1 & 0 & 0 \\
1 & 2 & 2 & 0 & 1 & 0 \\
1 & 3 & 2 & 0 & 0 & 1 \\
\end{block}
\end{blockarray}	\\
\xrightarrow{R_2 \rightarrow R_2-R_1} &
\begin{blockarray}{cccccc}
\begin{block}{[ccc|ccc]}
1 & 1 & 1 & 1 & 0 & 0 \\
0 & 1 & 1 & -1 & 1 & 0 \\
1 & 3 & 2 & 0 & 0 & 1 \\
\end{block}
\end{blockarray}	\\
\xrightarrow{R_3 \rightarrow R_3-R_1} &
\begin{blockarray}{cccccc}
\begin{block}{[ccc|ccc]}
1 & 1 & 1 & 1 & 0 & 0 \\
0 & 1 & 1 & -1 & 1 & 0 \\
0 & 2 & 1 & -1 & 0 & 1 \\
\end{block}
\end{blockarray}	\\
\xrightarrow{R_3 \rightarrow R_3-2R_2} &
\begin{blockarray}{cccccc}
\begin{block}{[ccc|ccc]}
\hlight{1} & 1 & 1 & 1 & 0 & 0 \\
0 & \hlight{1} & 1 & -1 & 1 & 0 \\
0 & 0 & \hlight{-1} & 1 & -2 & 1 \\
\end{block}
\end{blockarray}	\\
\xrightarrow{R_3 \rightarrow -1R_3} &
\begin{blockarray}{cccccc}
\begin{block}{[ccc|ccc]}
1 & 1 & 1 & 1 & 0 & 0 \\
0 & 1 & 1 & -1 & 1 & 0 \\
0 & 0 & 1 & -1 & 2 & -1 \\
\end{block}
\end{blockarray}	\\
\xrightarrow[R_1 \rightarrow R_1-R_3]{R_2 \rightarrow R_2-R_3} &
\begin{blockarray}{cccccc}
\begin{block}{[ccc|ccc]}
1 & 1 & 0 & 2 & -2 & 1 \\
0 & 1 & 0 & 0 & -1 & 1 \\
0 & 0 & 1 & -1 & 2 & -1 \\
\end{block}
\end{blockarray}	\\
\xrightarrow{R_1 \rightarrow R_1-R_2} &
\begin{blockarray}{cccccc}
\begin{block}{[ccc|ccc]}
\hlight{1} & 0 & 0 & 2 & -1 & 0 \\
0 & \hlight{1} & 0 & 0 & -1 & 1 \\
0 & 0 & \hlight{1} & -1 & 2 & -1 \\
\end{block}
\end{blockarray}.
\end{align*}
Det fremgår derfor, at $A$ er invertibel, da dens reducerede trappeform er $I_3$. $B$-siden, hvor $I_3$ før var, er nu
\begin{align*}
B=A^{-1}=\begin{blockarray}{ccc}
\begin{block}{[ccc]}
2 & -1 & 0 \\
0 & -1 & 1 \\
-1 & 2 & -1 \\
\end{block}
\end{blockarray}.
\end{align*}
%
\end{eks}
%%%%%%%%%%%%%%%%%%%%%%%%%%%%%%
Sætning \ref{thm:mxinvers2} beskriver nogle yderligere egenskaber for invertible matricer. 
%
\begin{thm}{}{mxinvers2}
%Givet to $m \times n$ matricer $A$ og $B$, så gælder, at (det er bare leftover aye?)
Lad $A$ være en $n \times n$ matrix. 
Da er følgende udsagn ækvivalente.
\begin{enumerate}[label=(\alph*)]
\item $A$ er invertibel.
\item Den reducerede trappeform af $A$ er $I_n$.
\item $\text{rang}(A)=n$.
\item Spannet af søjlerne i $A$ er $R^n$.
\item Ligningssystemet $A\textbf{x}=\textbf{b}$ er konsistent for alle $\textbf{b}$ i $R^n$.
\item $\text{null}(A)=0$.
\item Søjlerne i $A$ er lineært uafhængige.
\item Den eneste løsning for $A\textbf{x}=\textbf{0}$ er $\textbf{0}$.
\item Der eksisterer en $n \times n$ matrix $B$, så $BA=I_n$.
\item Der eksister en $n \times n$ matrix $C$, så $AC=I_n$.
\item $A$ er et produkt af elementære matricer.
%hvis vi tager det sidste punkt med er vi også nødt til at redegøre for det. 
\end{enumerate}
\end{thm}
\begin{proof}
Lad $A$ være en $n \times n$ matrix.
Jævnfør sætning \ref{thm:mxinvertibel} er (a) og (b) ækvivalente. 
Da $A$ er kvadratisk, er (b) ligeledes ækvivalent med (c), (d) og (e) jævnfør sætning \ref{thm:spaneqv}. 
Jævnfør sætning \ref{thm:mxlinuaf} er (b) desuden ækvivalent med (f), (g) og (h). 
Eftersom (a) er ækvivalent med (b), er (a) også ækvivalent med (c), (d), (e), (f), (g) og (h). 
\\\\
At (a) og (k) er ækvivalente, bevises i to skridt. 
Bemærk først, at hvis $A$ er invertibel, så er $I_n$ den reducerede trappeform af $A$. 
Jævnfør sætning \ref{thm:eleinv} og \ref{thm:mxinvertibel} eksisterer der en invertibel matrix $P$, så at $PA=I_n$, hvormed $A=P^{-1}$. 
Da $P=E_kE_{k-1} \cdots E_1$, så er $P^{-1}=E_1^{-1} \cdots E_{k-1}^{-1}E_k^{-1}$.
$A$ er dermed et produkt af elementære matricer, hvilket beviser, at (a) medfører (k). \\\\
Antag nu, at $A$ er et produkt af elementære matricer. 
Elementære matricer er invertible, hvormed $A$ er et produkt af invertible matricer. 
$A$ er derfor invertibel, hvilket bekræfter (a), hvilket gør (a) og (k) ækvivalente. 
\\\\
At (a) medfører (i) er åbenlyst, da $B=A^{1}$. 
Antag nu, at der eksisterer en $n \times n$ matrix $B$, så $BA=I_n$. Lad $\textbf{v} \in \R^n$, så at $A\textbf{v}=\textbf{b}$. Så gælder det, at 
$$
\textbf{v} = I_n\textbf{v} = (BA)\textbf{v} = B(A\textbf{v}) = B\textbf{0} = \textbf{0}.
$$
Bemærk, at (i) dermed medfører (h).
Eftersom (h) er ækvivalent med (a), medfører (i) derfor også (a).
Således er (i) og (a) ækvivalente. 
\\\\
At (a) medfører (j) er åbenlyst, da $C=A^{1}$.
Antag nu, at der eksister en $n \times n$ matrix $C$, så $AC=I_n$. 
Lad $\textbf{b} \in \R^n$ og $\textbf{v}=C\textbf{b}$. 
Så gælder, at 
$$
A\textbf{v} = A(C\textbf{b}) = (AC)\textbf{b} = I_n\textbf{b} = \textbf{b}.
$$
Bemærk, at (j) dermed medfører (e).
Eftersom (e) er ækvivalent med (a), medfører (j) derfor også (a).
Således er (j) og (a) ækvivalente, og alle punkter i sætning \ref{thm:mxinvers2} er dermed ækvivalente.
\end{proof}
	%
% EKSEMPEL! tilføjes 
%
\subsection{Invertibilitet af funktioner}
%
En to tresko - indlednende tekst.
%
\begin{defn}{}{}
En funktion $f:\R^n \rightarrow \R^m$ siges at være \textbf{surjektiv}, hvis dens værdimængde er hele $\R^m$, tilsvarende at alle vektorer i $\R^m$ er en afbildning af mindst én vektor i $\R^n$.
Funktionen siges at være \textbf{injektiv}, hvis hver entydig vektor i $\R^n$ har en entydig vektor i $\R^m$, som afbildning.
Hvis en funktion er både injektiv og surjektiv, så er funktionen \textbf{bijektiv}.
\end{defn}
\noindent
%
Bemærk, at
$$f^{-1} \neq \frac{1}{f},$$
men blot notation for den \textit{inverse funktion} til $f$.
Givet en funktion $T:\R^n \rightarrow \R^n$, som er en lineær afbildning med standardmatricen $A$, så er $T$ invertibel, hvis og kun hvis $A$ er invertibel. 
Lad A være en $n \times n$ invertibel matrix. 
Så haves for alle $\textbf{v}$, at
$$T_AT_{A^{-1}}(\textbf{v})=T_A(T_{A^{-1}}(\textbf{v}))=T_A(A^{-1}\textbf{v})=A(A^{-1}\textbf{v})=AA^{-1}\textbf{v}=I_n\textbf{v}=\textbf{v}.$$
Det kan gennem lignende operationer vises, at
$$T_{A^{-1}}T_A(\textbf{v})=\textbf{v}.$$
%
Dette medfører $T^{-1}=T_{A^{-1}}$. 
Som følge deraf er $T^{-1}$ lineær og dens standardmatrix er $A^{-1}$.
\\
%
\begin{eks}
Fra \ref{eks:fisk5} haves, at
\begin{equation*}
A=
\begin{bmatrix}
1 & 1 & 1\\
1 & 2 & 2\\
1 & 3 & 2
\end{bmatrix} \text{\phantom{jd}og\phantom{jd}}
A^{-1}=
\begin{bmatrix}
2 & -1 & 0 \\
0 & -1 & 1 \\
-1 & 2 & -1 \\
\end{bmatrix}
\end{equation*}  
så er
\begin{equation*}
T_A
\left(
\begin{bmatrix}
x_1\\
x_2\\
x_3
\end{bmatrix}
\right)
=
\begin{blockarray}{c}
\begin{block}{[c]}
x_1+x_2+x_3\\
x_1+2x_2+2x_3\\
x_1+3x_2+2x_3\\
\end{block}
\end{blockarray}.
\end{equation*}
Og den inverse funktion er
\begin{equation*}
T_{A^{-1}}\left(
\begin{bmatrix}
x_1\\
x_2\\
x_3
\end{bmatrix}
\right)=
\begin{blockarray}{c}
\begin{block}{[c]}
2x_1-1x_2\\
-x_2+x_3\\
-x_1+2x_2-x_3\\
\end{block}
\end{blockarray}.
\end{equation*}
\end{eks}
	% Underrum
	\section{Underrum}
Det kan ofte være belejligt at undersøge vektorer i en begrænset del af $\R^n$. Den begrænsede del kaldes et \textbf{underum}.
\begin{defn}{}{underrum}
En mængde $W$, bestående af vektorer i $\R^n$, er et underrum af $\R^n$, hvis den opfylder følgende betingelser:	
\begin{enumerate}[label=(\alph*)]
	\item $\textbf{0}\in W$.
	\item $\textbf{u}$,$\textbf{v} \in W \rightarrow \textbf{u}+\textbf{v} \in W $.
	\item $\textbf{u} \in W$ og $c$ er en skalar, så er $\textbf{u}c \in W$.
\end{enumerate}
\end{defn}
\noindent
Mængden $\R^n$ er et underrum til sig selv, da summen af to vilkårlige vektorer i $\R^n$ tilhører $\R^n$.
Enhver vektor multipliciteret med en skalar er i $\R^n$, og $\textbf{0}$ er ligeledes i $\R^n$. 
Mængden $\{\textbf{0}\}$ er ligeså et underrum til $\R^n$, kaldet \textbf{nul-underrummet}.
%
%
\begin{thm}{}{}
Spannet af en endelig ikke-tom delmængde af $\R^n$ er et underrum af $\R^n$.
\end{thm}
\begin{proof}
Lad spannet af en endelig ikke-tom mængde være 
$\S=\{\textbf{w}_1 , \textbf{w}_2 , \ldots , \textbf{w}_j\}$.
Eftersom der kan opskrives en linearkombination 
$$0\textbf{w}_1 + 0\textbf{w}_2 + \cdots + 0\textbf{w}_j=\textbf{0},$$ 
så er $\textbf{0}$ i  $\S$, hvilket opfylder betingelse (a) i definition \ref{defn:underrum}.
\\\\
%
Lad $\textbf{u}$ og $\textbf{v}$ tilhøre $\S$. Så er
%
\begin{align*}
\textbf{u}&=c_1\textbf{w}_1+c_2\textbf{w}_2+\cdots+c_j\textbf{w}_j, \\ 
\textbf{v}&=k_1\textbf{w}_1+k_2\textbf{w}_2+\cdots+k_j\textbf{w}_j,
\end{align*}
%
for skalarerne $c_1,c_2,\ldots,c_j$ og $k_1,k_2,\ldots,k_j$. 
Dette kan omskrives til linearkombinationen
%
\begin{align*}
\textbf{u}+\textbf{v}&=(c_1\textbf{w}_1+c_2\textbf{w}_2+\cdots+c_j\textbf{w}_j)+(k_1\textbf{w}_1+k_2\textbf{w}_2+\cdots+k_j\textbf{w}_j)\\
&=(c_1+k_1)\textbf{w}_1+(c_2+k_2)\textbf{w}_2+\cdots+(c_j+k_j)\textbf{w}_j,
\end{align*}
hvilket viser, at $\textbf{u}+\textbf{v}\in \S$, hvormed (b) i definition \ref{defn:underrum} er opfyldt.
\\\\
For enhver skalar $h$ gælder, at
\begin{align*}
h\textbf{u}&=h(c_1\textbf{w}_1+c_2\textbf{w}_2+\cdots+c_j\textbf{w}_j)\\
&=(h_1c_1)\textbf{w}_1+(h_2c_2)\textbf{w}_2+\cdots+(h_jc_j)\textbf{w}_j,
\end{align*}
og eftersom $h_1c_1,h_2c_2,\ldots,h_jc_j$ er skalarer, så er $h\textbf{u}\in \S$, hvilket opfylder (c) i definition \ref{defn:underrum}.
Hermed er det bevist, at spannet af en delmængde af $\R^n$ er et underrum til $\R^n$.
\end{proof}

\subsection{Underrum for matricer}
Der findes bestemte typer af underrum, som er relevante i forbindelse med matricer. Den første type, der beskrives, er \textbf{nulrummet}.
%
\begin{defn}{}{}
Nulrummet til en matrix $A$, noteret $\text{nrum}(A)$, er løsningsmængden til matrixligningen $A\textbf{x}=\textbf{0}$. 
\end{defn}
%
\noindent
Et eksempel på nulrummet til en matrix ses i eksempel \ref{nulrum}.

\begin{eks}
\label{nulrum}
Lad 
$$A=\begin{bmatrix}
2 & 3 & 2\\
1 & -1 & 5
\end{bmatrix}.$$
Matrixligningen $A\textbf{x}=\textbf{0}$,
kan opskrives som ligningssystemet
\begin{align*}
2x_1+3x_2+2x_3&=0\\
1x_1-x_2+5x_3&=0.
\end{align*}
Nulrummet, $\text{nrum}(A)$, til matricen er dermed de vektorer $\textbf{x}$, som opfylder ligningssystemet.
\end{eks} 
%
\begin{thm}{}{}
Hvis $A$ er en $m\times n$ matrix, så er nulrummet af $A$ et underrum af $\R^n$.
\end{thm}
%
\begin{proof}
Eftersom $A$ er en $m\times n$ matrix, så er vektorerne i nulrummet i $\R^n$, da det er løsningerne til $A\textbf{x}=\textbf{0}$. $\textbf{0}$ er en del af nulrummet, da $A\textbf{0}=\textbf{0}$. 
Hvis $\textbf{u}$ og $\textbf{v}$ er i nulrummet, så er $A\textbf{u}=\textbf{0}$ og $A\textbf{v}=\textbf{0}$. 
Deraf har vi, at
%
$$A(\textbf{u}+\textbf{v})=A\textbf{u}+A\textbf{v}=\textbf{0}+\textbf{0}=\textbf{0}.$$
%
Deraf haves, at $\textbf{u}+\textbf{v}$ er i nulrummet. 
For alle skalarer $c$ haves, at
$$A(c\textbf{u})=c(A\textbf{u})=c\textbf{0}=\textbf{0}.$$
Dette viser, at $c\textbf{u}$ er i nulrummet, hvilket viser, at nulrummet af $A$ er et underrum af $\R^n$
\end{proof}
\\
\noindent
Andre relevante typer underrum er \textbf{rækkerum} og \textbf{søjlerum}.
\begin{defn}{}{}
%Til nedenstående "Rækkerum er spannet af rækkerne i en matrix" - right? 
Rækkerum er spannet, som kommer af rækkerne i en matrix, og noteres $\text{ræk}(A)$.
\\
Søjlerum er spannet af søjlerne i en matrix, og noteres $\text{søj}(A)$.
\end{defn}
\noindent
Underrum med særegne karakteristika er også relevante i forbindelse med lineær afbildning.
Værdimængden af en lineær afbildning er det samme som søjlerummet af dens standardmatrix.
%
\subsection{Basis}
Mens et underrum kan beskrives med en given mængde vektorer, så er det fordelagtigt at beskrive det med så få vektorer som muligt. 
Denne mængde kaldes \textbf{basis} for underrummet.
%
\begin{defn}{}{basisunderrum}
Lad $\mathcal{P}$ være et ikke-nul-underrum af $\R^n$. 
En basis for $\mathcal{P}$ er en lineært uafhængig mængde af generatorer for $\mathcal{P}$.
\end{defn}
\noindent
Eksempelvis kan en mængde af standardvektorer $\{\textbf{e}_1,\textbf{e}_2,\ldots,\textbf{e}_n\}$ bruges som basis for $\R^n$.
Denne type basis kaldes for \textbf{standard basis}.
Der findes dog flere mulige baser for det samme underrum.
\\\\
%
Pivotsøjlerne i en matrix udgør basis for dens søjlerum. Dette følger af bevis \ref{thm:pivotu}.
\\\\
% Bevis mangler sætning 2.4 i bogen, bevis er i appendiks E xD
	\subsection{Dimension}
I forlængelse af underrum og basis introduceres nu
\textit{dimension}. 
%
\begin{defn}{}{dimension}
Mængden af vektorer, der udgør en basis for et givent ikke-nul-undderrum $V$ i $\R^n$, betegnes som dimensionen af $V$, hvilket noteres $\text{dim}(V)$. 
%Det er dertil fordelagtigt at definere dimensionen af nulrummet i $\R^n$ til at være 0. Nødvendigt?? NEEEEEEEEEEEEEEEEEEEEEEJ
\end{defn}
\noindent
Det betyder derfor generelt, at standardbasen for $\R^n$ indeholder $n$ vektorer, og derfor er $\text{dim}(\R^n) = n$. 
Dertil kan man betragte bestemte underrum associeret med matricer samt deres \textbf{dimensioner}. 
Der gælder for dimensionerne af underrummet for en $m \times n$ matrix, at pivotsøjlerne i enhver matrix danner en basis for dens søjlerum. 
Derfor er $\text{dim}(\text{søj}(A))$ lig med mængden af pivotsøjler i en matrix $A$ og dermed lig med $\text{rang}(A)$. 
Nulrummets dimension for en matrix er givet ved nulliteten af matricen $A$, altså $n-\text{rang}(A)$. 
\\\\
%
\begin{table}[h!]
\begin{center}
 \begin{tabular}{||c c c||} 
 \hline
 Underrum & Indeholder rum & Dimension\\
 \hline\hline
 søj$(A)$ & $\R^m$ & rang($A$)\\ 
 \hline
 nrum$(A)$ & $\R^n$ & null(A)= $n-\text{rang}(A)$\\
 \hline
 ræk$(A)$ & $\R^n$ & rang($A$)\\
 \hline
\end{tabular}
\caption{Oversigt over dimensioner af bestemte underrum for en $m \times n$ matrix.}
\end{center}
\end{table}
%Skal der være paranteser eller ej? Mads og jeg synes ja. (altså: Søj(A) eller Søj A?) Alle synes søj(A)
%


%
\noindent
Slutteligt er det gældende for rækkerummet for en matrix $A$, at dens dimension ligeledes er lig med matricens rang. 
Da rækkerummet og søjlerummet har samme dimension gælder, at
%
\begin{align*}
\text{dim}(\text{ræk}(A))=\text{dim}(\text{søj}(A))=\text{dim}(\text{ræk}(A^T)).
\end{align*}
%
Det følger derfor, at rangen af enhver matrix er lig med rangen af matricens transponerede. 
%
	

	
% LINEÆR OPTIMERINGSPROBLEMER 
	% Generel definition - Optimering 
	\chapter{Lineære optimeringsproblemer}
%
Lineær programmering er en problemstilling, hvor det ønskes at optimere et problem bestående af en lineær funktion og \textit{lineære betingelser}.
Funktionen, som optimeres, kaldes \textit{objektfunktionen}.
Objektfunktionen kan enten optimeres for maksimum eller minimum og er begrænset af de lineære betingelser.
Findes et maksimum eller minimum inden for begrænsningerne, kaldes dette for problemets \textit{optimale værdi}.
\\\\
%
Et lineært maksimumsproblem kan opstilles som følgende:
%
\begin{align*}
\begin{array}{lrll}
\text{Maksimér}		&\textbf{c}^T\textbf{x}	&			&\\
\text{begrænset af}	&\textbf{a}^T\textbf{x}	&\geq b_i,	&i \in M_1,\\
					&\textbf{a}^T\textbf{x}	&\leq b_i,	&i \in M_2,\\
					&\textbf{a}^T\textbf{x}	& = b_i,	&i \in M_3,\\
					&x_j					&\geq 0,	&j \in N_1,\\
					&x_j					&\leq 0,	&j \in N_2,
\end{array}
\end{align*}
%
hvor $\textbf{c}^T\textbf{x}$ er objektfunktionen og $\textbf{a}^T\textbf{x}$ er de lineære betingelser, samt $M_1, M_2, M_3, N_1, N_2$ er endelige mængder med indekser.
\textit{Løsningsmængden} betegner mængden af vektorer, som opfylder alle betingelserne.
Vektorer indeholdt i løsningsmængden kaldes for \textit{løsningsvektorer}.
Hvis værdiområdet er uendeligt stort, kaldes optimeringsproblemet for \textit{ubegrænset}; ellers kaldes det \textit{begrænset}, hvilket vil blive defineret senere i rapporten. 
\\\\
%
\begin{eks}
\label{eks:min_lin}
Betragt minimumsproblemet
%
\begin{align*}
\begin{array}{lrrlr}
\text{Minimér}		&	\multicolumn{2}{c}{4x_1-3x_2}  &\\
\text{begrænset af}	&2x_1& - x_2			&\leq 	&-10,\\
					&x_1& - x_2				&\geq	& 7,\\
					&x_1& + x_2				&\leq	& 20,\\
					&x_1&					&\geq	& 0,\\
					& &x_2					&\leq	& 12.
\end{array}
\end{align*}
%
På figur \ref{fig:min_beg} kan løsningsområdet, markeret med gråblå, for minimumsproblemet ses.
%
\begin{figure}[H]
\centering
\begin{tikzpicture}
% 
%
\begin{axis}[
	axis lines = left,
	xlabel = $x_1$,
	ylabel = \rotatebox{-90}{$x_2$},
	legend style={at={(1,0.05)},anchor=south east},
]
	\addplot[name path=b, domain={0:1}, color=myblue, style=very thick,samples=2]{2*x+10};
	\addlegendentry{$2x_1-x_2 \leq -10$}
	\addplot[name path=a, domain={0:13.5}, color=mygreen, style=very thick, samples=2]{x-7};
	\addlegendentry{$\phantom{.}x_1-x_2 \geq \phantom{....}7$}
	\addplot[domain=0:15,samples=2,color=myred, style=very thick]{-x+20};
	\addlegendentry{$\phantom{.}x_1 + x_2 \leq \phantom{..}20$}
	\addplot[name path=c, domain={1:8}, color=black, style=very thick,samples=2]{12};
	\addlegendentry{$\phantom{.........}x_2 \leq \phantom{..}12$}
	
	\tikzfillbetween [of=a and b]{myblue!15}
	\tikzfillbetween [of=c and a]{myblue!15}
	
	\addplot[domain=0:15,samples=2,color=myred, style=very thick]{-x+20};
	\addplot[domain=0:5,samples=2,color=myblue, style=very thick]{2*x+10};
	\addplot[domain=0:15,samples=2,color=mygreen, style=very thick]{x-7};
	\addplot[domain=0:15,samples=2,color=black, style=very thick]{12};
	
	\addplot[domain=0:15, samples=2, color=black, dotted, thick]{0};
	\addplot[domain=0:15, samples=2, color=black, dotted, thick]{0};
	\addplot[name path=c, domain=0:0.00003,samples=2,color=black, style=very thin]{10000*x};
	\addplot[mark=none] coordinates {(0,-7) (0,10)};
	
%
\end{axis}
\end{tikzpicture}
\caption{En begrænset løsningsmængde, markeret med blå, afgrænset af lineære betingelser.}
\label{fig:min_beg}
\end{figure}
\noindent
%
Med udgangspunkt i figur \ref{fig:min_beg} kan det ses, at løsningsmængden er begrænset.
% 
Havde minimumsproblemet i stedet været 
%
\begin{align*}
\begin{array}{lrrlr}
\text{Minimer}		&\multicolumn{2}{c}{4x_1-3x_2}  &\\
\text{begrænset af}	&2x_1& - x_2			&\leq 	&-10,\\
					&x_1& - x_2				&\geq	& 7,\\
					&x_1&					&\geq	& 0,
\end{array}
\end{align*}
havde minimumsproblemet været ubegrænset, da $x_1$ og $x_2$ kunne blive uendeligt store. 
Et undsit af løsningsområdet for minimumsproblemet ses på figur \ref{fig:min_ubeg}, hvor det er markeret med blå.
%
\begin{figure}[H]
\centering
\begin{tikzpicture}
\begin{axis}[
	axis lines = left,
	xlabel = $x_1$,
	ylabel = {$x_2$},
	legend style={at={(0.45,0.025)},anchor=south west},
]
	\addplot[name path=b, domain={0:5}, color=myblue, style=very thick,samples=2]{2*x+10};
	\addlegendentry{$2x_1-x_2 \leq -10$}
	\addplot[name path=a, domain={0:15}, color=mygreen, style=very thick, samples=2]{x-7};
	\addlegendentry{$x_1-x_2 \geq 7$}
	\addplot[name path=c, domain={5:15}, color=white, style=very thick,samples=2]{20};
	
	\tikzfillbetween [of=a and b]{myblue!15}
	\tikzfillbetween [of=c and a]{myblue!15}
	
	\addplot[domain=0:5,samples=2,color=myblue, style=very thick]{2*x+10};
	\addplot[domain=0:15,samples=2,color=mygreen, style=very thick]{x-7};
	
	\addplot[domain=0:15, samples=2, color=black, dotted, thick]{0};

\end{axis}
\end{tikzpicture}
\caption{Betingelserne, der danner ubegrænset løsningsmængden, markeret med blå, for minimeringsproblemet.}
\label{fig:min_ubeg}
\end{figure}

%
\end{eks}
	% Standartform 
	\section{Standardform}
% 

\begin{align*}
\begin{array}{lrl}
\text{Maksimer}		&\textbf{c}^T\textbf{x}	&				\\
\text{begrænset af}	&\textbf{A}\textbf{x}	&=\mathbf{b},	\\
					&\mathbf{x}				&\geq \mathbf{0},		
\end{array}
\end{align*}
	% Geometrisk løsning 
	\section{Mathias}
% God arbejdslyst - tak julle :))

I forhold til grafisk at løse lignende lineære programmeringsproblemer forhenstående, så vil følgende eksempel illustrere hvordan. 
%Mit forslag til ovenstående: Eksempel \ref{eks:min_loes} illustrerer, hvordan lineære programmeringsproblemer løses i nærværende rapport.

\begin{eks}
\label{eks:min_loes}
Betragt minimumsproblemet
%
\begin{align*}
\begin{array}{lrrlr}
\text{Minimér}		&	\multicolumn{2}{c}{x_1-x_2}  &\\
\text{begrænset af}	&x_1& + 2x_2			&\leq 	& 6,\\
					&2x_1& + x_2			&\leq	& 6,\\
					&x_1&    				&\geq	& 0,\\
					&x_2&    				&\geq	& 0.\\
\end{array}
\end{align*}
%
Grafisk illustreret ved 
%
\begin{figure}[H]
\centering
\begin{tikzpicture}
\begin{axis}[
	 axis x line=center,
  axis y line=center,
  xtick={-5,-4,...,5},
  ytick={-5,-4,...,5},
  xlabel={$x_1$},
  ylabel={$x_2$},
  xlabel style={below right},
  ylabel style={above left},
  xmin=-1,
  xmax=4,
  ymin=-1,
  ymax=4]
	legend style={at={(0.45,0.025)},anchor=south west},
]



\addplot[name path=b, domain={0:6}, color=myblue, style=very thick,samples=2]{6-2*x};
	\addlegendentry{$2x_1+x_2 \leq 6$}
	\addplot[name path=a, domain={0:6}, color=myred, style=very thick, samples=2]{3-0.5*x};
	\addlegendentry{$x_1+2x_2 \leq 6$}
	

	\addplot[domain=0:6,samples=2,color=myred, style=very thick]{((6-2*x};
	\addplot[domain=0:6,samples=2,color=myblue, style=very thick]{3-0.5*x};

	
\addplot[domain=-1:6, samples=2, color=black, dotted, thick]{-x};
\addplot[domain=-1:6, samples=2, color=black, dotted, thick]{2-x};
\addplot[domain=0:6, samples=2, color=black, dotted, thick]{4-x};
\addplot[mark=*] coordinates {(2,2)};
\addplot[->] coordinates {(0,0)--(-1,-1)};


\end{axis}
\end{tikzpicture}
\caption{Grafisk løsning til problemet i \ref{eks:min_loes}}
\label{fig:min_loes}
\end{figure}
\noindent
Løsningsmængden er det skraverede område på figur \ref{fig:min_loes}. For at finde den optimale løsning findes enhver given skalar $k$ og overvejer løsningsmængden med alle punkter, hvor objektfunktionen $\textbf{c}^T\textbf{x}$ er lig med $k$. Dette illustreres ved den linje, der er beskrevet af ligningen $-x_l-x_2=k$. Bemærk, at denne linje er vinkelret på vektoren $\textbf{v}$ med koordinaterne $\textbf{v}=(-1,-1)$.
Forskellige værdier for $k$ fører til forskellige linjer, hvor  alle linjer er parallelle med hinanden samt vinkelret med vektoren. Hvis man øger værdien af $k$ svarer til at bevæge linjen $z=x_l-x_2$ langs retningen af vektoren $\textbf{v}$. For at finde den optimale løsning er det interessante at minimere k, og dermed bevæge linjen så meget som muligt i retning af $-\textbf{v}$ uden at forlade løsningsmængden. Det mest optimale er derfor $k = -4$ og dette gør, at vektor $\textbf{x}=(2,2)$ er den mest optimale løsning i løsningsmængden. 
\end{eks}
\noindent
Løsningsmængden er det skraverede område på figur \ref{fig:min_loes}. For at finde den optimale løsning findes enhver given skalar $k$ og løsningsmængden med alle punkter overvejes, hvor objektfunktionen $\textbf{c}^T\textbf{x}$ er lig med $k$. Dette illustreres ved den linje, der er beskrevet af ligningen $-x_l-x_2=k$. Dette betyder, at den optimale løsning som enten maksimerer eller minimerer værdien
af objektfunktionen kan findes grafisk ved at indtegne niveaukurver $f(x1, x2)=k$ og forskyde disse niveaukurver i den retning, der optimerer værdien af objektfunktionen, indtil den optimale løsning opnås.  Bemærk, i det følgende eksempel, at denne linje er vinkelret på vektoren $\textbf{v}$ med koordinaterne $\textbf{v}=(1,1)$.
Forskellige værdier for $k$ fører til forskellige linjer, hvor  alle linjer er parallelle med hinanden samt vinkelret med vektoren $\textbf{v}$. Hvis værdien af $k$ øges, svarer det til at bevæge linjen $z=x_l-x_2$ langs retningen af vektoren $\textbf{v}$. For at finde den optimale løsning for dette eksempel, er det interessante at maksimere $k$, og dermed bevæge linjen så meget som muligt i retning af $\textbf{v}$ uden at forlade løsningsmængden. Det mest optimale er derfor $k = -4$, hvormed vektor $\textbf{x}=(2,2)$ er den mest optimale løsning i løsningsmængden. \\\\
%%%%%
Det hjørnepunkt, som er den optimale løsning i forudgående eksempel, som i forudgående tilfælde maksimere værdien
af objektfunktionen repræsenterer præcis én
løsning til maksimeringsproblemet. I følgende eksempel introduceres et eksempel, hvor der matematisk set findes et uendeligt antal løsninger, der
maksimerer værdien af objektfunktionen.

	% Dualitet 
	\section{Dualitet}
% Vi kan vælge at tage det med lagrange multiplier med også men tænker vi snakker om det jeg er personligt i tvivl og ville nødigt spilde tiden på det.
	% - Julie: Jeg synes ikke vi skal have det med 
% Jeg er iøvrigt bange for jeg er kommet til at skyde mig selv i foden, men umiddelbart giver det mening for mig, det snakker vi om.
	% - Julie: Jeg tror du har styr på det xD
%
Et vigtigt element af teorien vedrørende lineære optimerings problemer og Simplex metodens brug, omhandler \textit{dualproblemer}.
%
\begin{defn}{}{dualproblemer}
Givet maksimumsproblemet: 
%
\begin{align*}
\begin{array}{lrl}
\text{Maksimér}		&\textbf{c}^T\textbf{x}	& = z		\\
\text{begrænset af}	&A\textbf{x}	&\leq \mathbf{b},	\\
					&\mathbf{x}				&\geq \mathbf{0},
\end{array}
\end{align*}
%
Da vil \textbf{dual minimumsproblemet} være: \\
% Minimer $v=\mathbf{b}^T \mathbf{p}$ begrænset af $A^T \mathbf{p} \geq c$ og $\mathbf{p} \geq 0$
%
\begin{align*}
\begin{array}{lrl}
\text{Minimer}		&\textbf{b}^T\textbf{p}	& = v			\\
\text{begrænset af}	&A^T \textbf{p}	&\leq \mathbf{c},	\\
					&\mathbf{p}				&\geq \mathbf{0},
\end{array}
\end{align*}
%
\end{defn}
%
Problemet:
%
\begin{align*}
\begin{array}{lrll}
\text{Maksimér}		&\textbf{c}^T\textbf{x}	&			&\\
\text{begrænset af}	&\textbf{a}_i^T\textbf{x}	&\geq b_i,	&i \in M_1,\\
					&\textbf{a}_i^T\textbf{x}	&\leq b_i,	&i \in M_2,\\
					&\textbf{a}_i^T\textbf{x}	& = b_i,	&i \in M_3,\\
					&x_j					&\geq 0,	&j \in N_1,\\
					&x_j					&\leq 0,	&j \in N_2,\\							&x_j					&\text{fri},	&j \in N_3.
\end{array}
\end{align*}

Vil således have:


\begin{align*}
\begin{array}{lrll}
\text{minimer}		&\textbf{p}^T\textbf{b}	&			&\\
\text{begrænset af}	&p_i					&\geq 0,	&i \in M_1,\\
					&p_i					&\leq 0,	&i \in M_2,\\
					&p_i					&\text{fri},	&i \in M_3,\\
					&\textbf{p}^T\textbf{A}_j	&\leq c_j,	&j \in N_1,\\
					&\textbf{p}^T\textbf{A}_j	&\geq c_j,	&j \in N_2,\\
					&\textbf{p}^T\textbf{A}_j	& = b_i,	&i \in N_3,
\end{array}
\end{align*}
%
som sit dualproblem.
%Ideer til en federe opsætning
\\\\
%
Det sker således at der for hver begrænsning i det oprindelige problem indføres en variable i dualproblemet, tilsvarende indføres der en begrænsning for hver variable. 
Derudover skiftes som nævnt i definitionen fra et maksimerings til minimeringsproblem. 
Dette kan opsummeres i følgende tabel: \\
\begin{table}[H]
\begin{center}
\begin{tabular}{llll}
Oprindelige  & maksimer   & minimer    & Dual         \\
\hline
Variable     & $\geq 0$   & $\geq b_i$ & Begræsninger \\
             & $\leq 0$   & $\leq b_i$ &              \\
             & fri        & $=b_i$     &              \\
\hline
Begræsninger & $\leq c_j$ & $\geq 0$   & Variable     \\
             & $\geq c_j$ & $\leq 0$   &              \\
             & $=c_j$     & fri        &             
\end{tabular}
\end{center}
%\captionof{tabular}{Fisk finder i på noget.}
% Der mangler noget tabeltekst ved ikke lige hvordan spørg MADS xD
\end{table}
\noindent
%
Dual problemet kan således opstilles som i eksempel \ref{dual} \\
%
\begin{eks}
\label{dual}
%
Betragt diæt problemet fra afsnittet om simplexmetoden:
%
\begin{align*}
\begin{array}{lrrlr}
\text{Maksimer}		&	\multicolumn{2}{c}{z=5x_1+4x_2}  &\\
\text{begrænset af}	&3x_1& +5x_2			&\leq 	&78,\\
					&4x_1& + x_2				&\leq	& 36,\\
					&x_1& , x_2				&\geq	& 0.
\end{array}
\end{align*}
\\
Dualproblemet vil således have denne form:
%hvilket bogstave foretrækker i istedet for z? hvis i vil ha z i standard notationen 
	% Julie: Tænker v? var det ikke det du skrev i den forrige?
\begin{align*}
\begin{array}{lrrlr}
\text{Minimer}		&	\multicolumn{2}{c}{z=78p_1+36_p2}  &\\
\text{begrænset af}	&3p_1& +4p_2			&\geq 	&5,\\
					&5p_1& + p_2				&\geq	& 4,\\
					&p_1& , p_2				&\geq	& 0.
\end{array}
\end{align*}
Hvor problemet her istedet er at minimere indholdet af sukker og kulhydrater begrænset af forholdet mellem to typer slik i en slikpose.
%?? Julie skal have lidt mere forklaring af dine begrænsninger for at komme i mål her.
%
Løses dette ved hjælp af simplex-algoritmen opnås følgende tabel:
\begin{align*}
\begin{blockarray}{cccccc}
p_1 & p_2 & \textcolor{blue}{s_1} & \textcolor{blue}{s_2} & z & b \\
\begin{block}{[cc|ccc|c]}
3 & 4 & 1 & 0 & 0 & 5 \\
5 & 1 & 0 & 1 & 0 & 4 \\
\cline{1-6}
78 & 36 & 0 & 0 & 1 & 0\\
\end{block}
\end{blockarray}
\end{align*}
%skal have lidt hjælp til at omskrive problemet til standardform simplex(det er bare fortegn jeg er i tvivl om, har et program der udføre algoritmen, så er good to go bagefter)
%
%her skal løst simplex smides ind
%
Som det fremgår findes minimum i samme punkt:
$z=78$ hvor $p_1$ og $p_2$ slackvariable $s_1=6$ og $s_2=12$
%
\end{eks}
Som det fremgår i eksemplet har løsningen på dualproblemet samme optimale løsning, hvor slackvariable her er løsningen på det oprindelige problem.
Det gælder endvidere at der fra dualproblemet kan vendes tilbage til det oprindelige problem hvilket udtrykkes ved følgende sætning:
\begin{thm}{}{dualitetsaet}
Hvis dualproblemet omdannes til et ækvivalent minimumsproblem og dualproblemet af dette dannes, opnås det oprindelige problem.
\end{thm}
\begin{proof}
Betragt problemet:
\begin{align*}
\begin{array}{lrll}
\text{minimer}		&\textbf{p}^T\textbf{b}	&			&\\
\text{begrænset af}	&p_i					&\geq 0,	&i \in M_1,\\
					&p_i					&\leq 0,	&i \in M_2,\\
					&p_i					&\text{fri},	&i \in M_3,\\
					&\textbf{p}^T\textbf{A}_j	&\leq c_j,	&j \in N_1,\\
					&\textbf{p}^T\textbf{A}_j	&\geq c_j,	&j \in N_2,\\
					&\textbf{p}^T\textbf{A}_j	& = b_i,	&i \in N_3,
\end{array}
\end{align*}
og dets dualproblem:
\begin{align*}
\begin{array}{lrll}
\text{Maksimér}		&\textbf{c}^T\textbf{x}	&			&\\
\text{begrænset af}	&\textbf{a}_i^T\textbf{x}	&\geq b_i,	&i \in M_1,\\
					&\textbf{a}_i^T\textbf{x}	&\leq b_i,	&i \in M_2,\\
					&\textbf{a}_i^T\textbf{x}	& = b_i,	&i \in M_3,\\
					&x_j					&\geq 0,	&j \in N_1,\\
					&x_j					&\leq 0,	&j \in N_2,\\							&x_j					&\text{fri},	&j \in N_3,
\end{array}
\end{align*}
Dualproblemet omdannes til et ækvivalent minimumsproblem:
\begin{align*}
\begin{array}{lrll}
\text{minimer}		&-\textbf{c}^T\textbf{x}	&			&\\
\text{begrænset af}	&-\textbf{a}_i^T\textbf{x}	&\leq -b_i,	&i \in M_1,\\
					&-\textbf{a}_i^T\textbf{x}	&\geq -b_i,	&i \in M_2,\\
					&-\textbf{a}_i^T\textbf{x}	& = -b_i,	&i \in M_3,\\
					&x_j					&\geq 0,	&j \in N_1,\\
					&x_j					&\leq 0,	&j \in N_2,\\							&x_j					&\text{fri},	&j \in N_3,
\end{array}
\end{align*}
Dualen til dette problem findes nu:
\begin{align*}
\begin{array}{lrll}
\text{maksimer}		&-\textbf{p}^T\textbf{b}	&			&\\
\text{begrænset af}	&p_i					&\geq 0,	&i \in M_1,\\
					&p_i					&\leq 0,	&i \in M_2,\\
					&p_i					&\text{fri},	&i \in M_3,\\
					&-\textbf{p}^T\textbf{A}_j	&\geq -c_j,	&j \in N_1,\\
					&-\textbf{p}^T\textbf{A}_j	&\leq -c_j,	&j \in N_2,\\
					&-\textbf{p}^T\textbf{A}_j	& = -b_i,	&i \in N_3,
\end{array}
\end{align*}
Hvilket har et ækvivalent minimumsproblem:
\begin{align*}
\begin{array}{lrll}
\text{minimer}		&\textbf{p}^T\textbf{b}	&			&\\
\text{begrænset af}	&p_i					&\geq 0,	&i \in M_1,\\
					&p_i					&\leq 0,	&i \in M_2,\\
					&p_i					&\text{fri},	&i \in M_3,\\
					&\textbf{p}^T\textbf{A}_j	&\leq c_j,	&j \in N_1,\\
					&\textbf{p}^T\textbf{A}_j	&\geq c_j,	&j \in N_2,\\
					&\textbf{p}^T\textbf{A}_j	& = b_i,	&i \in N_3,
\end{array}
\end{align*}
\end{proof}
%
%42 min https://youtu.be/0TRxEvMRE7s definition
	
	
% GEOMETRISK FREMSTILLING 
	\chapter{Grafisk løsning}
% 2.1 Polyeder og konvekse mængder 
% ---------------------------------------------------------
	\input{incl/main/grafisk_losning/hyperplan_halvrum_polyeder}
	\input{incl/main/grafisk_losning/konvekse_maengder}
% 
% 2.2 Ekstrema, hjørnepunkter og basale mulige løsninger 
% ---------------------------------------------------------
	\section{Ekstremer, hjørnepunkter og basale løsninger}
%
%\textit{Vis at løsningerne findes i hjørnene af en polytede.}
Som det fremgår af afsnit \ref{heeeeejjulle}, vil en optimal løsning være tilbøjelige til at ligge i hjørnerne af en polyede.
Der er derfor behov for måder at definere disse hjørner på.
%
\begin{defn}{}{ekstrema}
Lad $\mathcal{P}$ være en polyede. 
En vektor $\mathbf{x} \in \mathcal{P}$ kaldes et \textbf{ekstremumspunkt} i $\mathcal{P}$, hvis der ikke eksisterer to vektorer $\mathbf{y},\mathbf{z} \in \mathcal{P}$, $\mathbf{y} \land \mathbf{z} \neq \mathbf{x}$, 
samt en skalar $\lambda \in [0,1]$, hvorom det gælder: $\mathbf{x}=\lambda\mathbf{y}+(1-\lambda)\textbf{z}$.
\end{defn}
\noindent
%
%
Et hjørne kan beskrives ud fra denne definition idet der såfremt $\mathbf{x}=\lambda\mathbf{y}+(1-\lambda)z$, og alle vektorene findes i $P$, gælder at $\mathbf{x}$ er en konveks kombination af $\mathbf{y}$ og $\mathbf{z}$.
Hvis $\mathbf{x}=\lambda\mathbf{y}+(1-\lambda)\textbf{z}$ og $\mathbf{x}$ er et ekstremumspunkt må det derfor gælde at $\mathbf{y}\notin P$ eller $\mathbf{z}\notin P$ eller $\mathbf{x}=\mathbf{z}$ eller $\mathbf{x}=\mathbf{y}$. 
Det skal her nævnes at denne definition er strengt geometrisk.
%

\begin{figure}[h!]
  \centering
  \begin{tikzpicture}
    \tikzset{punkt/.style={point, draw=black}}\draw[color=black](1,-2) circle (2.5);
	\draw[color=black](7,-2) circle (2.5);  
	\draw[color=black, fill=myblue!15](7,-2) circle (1.5);  
%Punkter
	\draw  node[fill,circle,inner sep=0pt,minimum size=3pt] at (1,-2.65)  (v1) {};
	\draw  node[fill,circle,inner sep=0pt,minimum size=3pt] at (1,-1.35)  (v2) {};
	\draw  node[fill,circle,inner sep=0pt,minimum size=3pt] at (7,-2)  (v3) {};

%Labels
	\draw  node at (0.75,-2.55)  {$\textbf{u}$};
	\draw  node at (0.75,-1.25)  {$\textbf{v}$};
	\draw  node at (7.25,-1.85)  {$\textbf{w}$};
	
	
	
	%Navn på cirklerne 
	\draw[black, fill=black] (1.25,1) circle (0pt) node[anchor=west] {$\S_1$};    
	\draw[black, fill=black] (6.75,1) circle (0pt) node[anchor=west] {$\S_2$};
	%Navn på punkter 

	\draw[black, fill=black] (4,1.5) circle (0pt) node[anchor=west] {$f$};
	
	\draw [->,thick, draw=black] (2,1) -- (6.5,1);

	 \draw [->, thick, draw=red] ([xshift=5pt, yshift=-1pt]v1.north) -- ([xshift=-5pt, yshift=1pt]v3.south);
     \draw [->, thick, draw=black] ([xshift=5pt, yshift=1pt]v2.south) -- ([xshift=-5pt, yshift=-1pt]v3.north);


	\draw[black, fill=black] (0,-5) circle (0pt) node[anchor=west] {Domæne};
	\draw[black, fill=black] (6,-5) circle (0pt) node[anchor=west] {Codomæne};
	\draw[black, fill=black] (5.5,-2.5) circle (0pt) node[anchor=west] {Værdimængde};    
	%Streger mellem punkterne
  \end{tikzpicture}
  \caption{En funktion fra $\S_1$ til $\S_2$, hvor $f(\textbf{v})=\textbf{w}$ og f$(\textbf{u})=\textbf{w}$. }
  \label{fig:afbild}
\end{figure}
%
%
%    \node[punkt] at (-4,0.5)      (v1){$v_1$};
%    \node[punkt] at (-2,0.5)      (v2){$v_2$};
%    \node[punkt] at (-4,-1.5)     (v3){$v_3$};
%    \node[punkt] at (-2,-1.5)     (v4){$v_4$};
%    \node at (-3,2)     (v){$K_{4}$};
%
%
%    \node[punkt] at (4.6,-0.2)      (k1){$v_3$};
%    \node[punkt] at (1.4,-0.2)      (k2){$v_2$};
%    \node[punkt] at (2,-2)     (k3){$v_4$};
%    \node[punkt] at (4,-2)     (k4){$v_5$};
%    \node[punkt] at (3,1)      (k5){$v_1$};
%    \node at (3,2)      (k){$K_{5}$};
%
%
%
%
%
%    \draw [-, thick, draw=black] (v1) -- (v2);
%    \draw [-, thick, draw=black] (v1) -- (v3);
%    \draw [-, thick, draw=black] (v1) -- (v4);
%    \draw [-, thick, draw=black] (v2) -- (v3);
%    \draw [-, thick, draw=black] (v2) -- (v4);
%    \draw [-, thick, draw=black] (v3) -- (v4);
\\\\
%
En alternativ geometrisk definition relatere sig til \textit{hjørne punkter}, som er den entydige optimale løsning til et givet lineært programmeringsproblem med den mulige løsningsmængde $\mathcal{P}$.
%
\begin{defn}{}{hjoerner}
Lad $P$ være en polyede. 
En vektor $\mathbf{x}\in P$ siges at være et \textbf{hjørne punkt}, hvis der esksiterer en vektor $\mathbf{c}$, hvorom det gælder at $\mathbf{c}^T\mathbf{x}<\mathbf{c}^T\mathbf{y}$ for alle $\mathbf{y}$, som opfylder $\mathbf{y} \in P$ samt $\mathbf{y}\neq\mathbf{x}$.
\end{defn}
\noindent
%
%
Dette kan ligeledes beskrives som at $\mathbf{x}$ er et hjørne i $P$ såfremt $P$ er på den ene side af et hyperplan der skærer $P$ i $\mathbf{x}$. 
Jævnfør definitionen har dette hyperplan ligningen $y \mid \mathbf{c}^T\mathbf{x}=\mathbf{c}^T\mathbf{y}$	
	%
\begin{defn}{}{algh}
Hvis en vektor $\textbf{x}$ opfylder $\textbf{a}^T_i\textbf{x}= b_i$ for $i \in M_1, M_2 \text{ eller } M_3$, så siges den tilsvarende betingelse at være \textbf{bindende} eller \textbf{aktive} ved $\textbf{x}$.
\end{defn}\noindent
%
Følgende sætning giver anledning til, at disse bindende betingelser kan føres i relation til løsninger af lineære optimeringsproblemer.
%
\begin{thm}{}{betingi}

Lad $\textbf{x}$ være et element i $\R^n$ og lad $I=\{i \text{ } | \text{ } \textbf{a}^T_i\textbf{x}=b_i\}$ være en mængde af indekser på betingelser, der er bindende ved $\textbf{x}$.
Så er følgende udsagn ækvivalente:
%
\begin{enumerate}[label=(\alph*)]
\item Der eksisterer $n$ vektorer i mængden $\{\textbf{a}_i \text{ } | \text{ } i \in I \}$, som er lineært uafhængige.
\item Spannet af vektorerne $\textbf{a}_i$, for $i \in I$, dækker hele $\R^n$.
\item Ligningssystemet $\textbf{a}^T_i\textbf{x}= b_i$, for $i \in I$, har en entydig løsning.
\end{enumerate}
\end{thm}
%
\begin{proof}
Antag, at $\text{span}\{a_i \text{ } | \text{ } i \in I\}= \R^n$.
Dette span har dimensionen $n$ og jævnfør sætning \ref{} %(1.3(a) i bertsimas) 
kan en basis for spannet dannes med $n$ af vektorerne, som vil være lineært uafhængige.
%
Hvis der ligeledes eksisterer $n$ vektorer i $\{a_i|i \in I\}$, som er lineært uafhængige, må disse udspænde hele $\R^n$.
Dette viser, at (a) og (b) er ækvivalente.\\\\
%
Antag at ligningssystemet $\textbf{a}^T_i\textbf{x}=b_i$, $i \in I$, har løsningerne $\textbf{x}_1$ og $\textbf{x}_2$.
Så skal ikke-nulvektoren $\textbf{d} = \textbf{x}_1 - \textbf{x}_2$ opfylde, at $\textbf{a}^T_i\textbf{d}=0$ for alle $i \in I$.
Dette kræver, at $\textbf{d}$ er ortogonal med alle $\textbf{a}_i$, $i \in I$ og dermed ikke er en linearkombination af disse vektorer.
Dette medfører, at $\text{span}\{a_i \text{ } | \text{ } i \in I\} \neq \R^n$.
%
Antag nu, at vektorerne $\textbf{a}_i$, $i \in I$ ikke spænder over $\R^n$.
Så er det muligt at vælge en ikke-nulvektor $\textbf{d}$, der er ortogonal med underrummet, som vektorerne spænder over.
Hvis $\textbf{x}$ opfylder  $\textbf{a}^T_i\textbf{x}= b_i$, for alle $i \in I$, så haves, at $\textbf{a}^T_i(\textbf{x}+\textbf{d})= b_i$, for alle $i \in I$, og der er dermed adskillige løsninger.
Dermed er det vist, at (b) og (c) er ækvivalente.
\end{proof}\\
%
Med udgangspunkt i betingelsers tilsvarende vektorer kan givne betingelser siges at være lineært uafhængige, hvis de tilsvarende vektorer er lineært uafhængige.
Med dette og sætning \ref{thm:betingi}(a) kan hjørnepunkterne defineres som mulige løsninger, hvor $n$ lineært uafhængige betingelser er bindende.
Bemærk, at alle betingelser skal være opfyldt, for at et givent punkt kan være en mulig løsning, samt at et punkt ikke behøver at opfylde alle betingelser for at kunne have $n$ bindende betingelser.
%
\begin{defn}{}{basal}
Lad $\mathcal{P}$ være en polyede defineret med lineære ligheds- og ulighedsbetingelser og lad $\textbf{x}$ være en vektor i $\R^n$.
%
\begin{enumerate}[label=(\alph*)]
\item Vektoren $\textbf{x}$ er en \textbf{basal løsning}, hvis:
%
\begin{enumerate}[label=(\roman*)]
\item Alle lighedsbetingelser er opfyldt.
\item Af de bindende betingelser er der $n$ af dem, som er lineært uafhængige.
\end{enumerate}
%
\item Vektoren $\textbf{x}$ er en \textbf{basal mulig løsning}, hvis $\textbf{x}$ er en \textit{basal løsning}, og alle betingelser er opfyldt.
\end{enumerate}
\end{defn}\noindent
%
Bemærk, at hvis en polyede er defineret med $m$ lineære betingelser og $m<n$, så eksisterer der ikke et punkt med $n$ bindende betingelser.
Dette medfører, at der hverken er basale løsninger eller basale mulige løsninger.\\\\
%
Tre definitioner, som ønsker at beskrive det samme, er blevet givet.
Det er muligt at skifte mellem definitionerne efter ønske, da de er ækvivalente.
%
\begin{thm}{}{hjorneeq}
Lad $\mathcal{P}$ være en ikke-tom polyede og lad vektoren $\textbf{x}\in \mathcal{P}$.
Så er følgende udsagn ækvivalente:
%
\begin{enumerate}[label=(\alph*)]
\item $\textbf{x}$ er et \textit{ekstremumpunkt}.
\item $\textbf{x}$ er et \textit{hjørnepunkt}.
\item $\textbf{x}$ er en \textit{basal mulig løsning}.
\end{enumerate}
%
\end{thm}
%
\begin{proof}
Lad $P$ være en polyede defineret ved betingelser på formen $\textbf{a}_i^T\textbf{x} \geq b_i$ og $\textbf{a}_i^T\textbf{x} = b_i$.\\\\
%
Lad $\textbf{x} \in \mathcal{P}$ være et hjørnepunkt. Jævnfør definition \ref{defn:hjoerner} eksisterer der så et $\textbf{c} \in \R^n$, således at $\textbf{c}^T\textbf{x} < \textbf{c}^T\textbf{y}$ for alle $\textbf{y}$, som opfylder $\textbf{y} \in \mathcal{P}$ og $\textbf{y} \neq \textbf{x}$.
Hvis $\textbf{y} \in \mathcal{P}$, $\textbf{z} \in \mathcal{P}$, $\textbf{y} \neq \textbf{x}$, $\textbf{z} \neq \textbf{x}$ og $0 \leq \lambda \leq 1$, så $\textbf{c}^T\textbf{x} < \textbf{c}^T\textbf{y}$ og $\textbf{c}^T\textbf{x} < \textbf{c}^T\textbf{z}$.
Dette medfører, at $\textbf{c}^T\textbf{x} < \textbf{c}^T(\lambda \textbf{y} + (1-\lambda)\textbf{z})$, som betyder, at $\textbf{x}^T \neq \lambda \textbf{y} + (1 - \lambda)\textbf{z}$.
Derfor kan $\textbf{x}$ ikke blive udtrykt som en kombination af to andre elementer i $\mathcal{P}$, hvilket betyder, at $\textbf{x}$ er et ekstremumspunkt ifølge definition \ref{defn:ekstrema}.
Hermed (b) $\Rightarrow$ (a).
\\\\
%
Antag, at $\textbf{x} \in \mathcal{P}$ ikke er en basal mulig løsning og lad $I = \{ i \text{ } | \text{ }  \textbf{a}_i^T \textbf{x} = b_i\}$.
Da $\textbf{x}$ ikke er en basal mulig løsning, må der ikke eksistere $n$ lineært uafhængige vektorer af typen $\textbf{a}_i$, $i \in I$.
Derfor må vektorerne $\textbf{a}_i$, $i \in I$, ligge i et ægte underrum af $\R^n$ og der eksisterer derfor en ikke-nulvektor $\textbf{d} \in \R^n$, således at $\textbf{a}_i^T\textbf{d} = 0$ for alle $i \in I$.
Lad $\epsilon$ være et lille, positivt tal, og betragt vektorerne $\textbf{y} = \textbf{x} + \epsilon\textbf{d}$ og $\textbf{y} = \textbf{x} - \epsilon\textbf{d}$.
Bemærk, at $\textbf{a}_i^T \textbf{y} = \textbf{a}_i^T \textbf{x} = b_i$, for $i \in I$.
Endvidere haves, at $\textbf{a}_i^T \textbf{x} > b_i$ for $i \notin I$ og givet, at $\epsilon$ er lille nok, haves $\textbf{a}_i^T \textbf{y} > b_i$ for $i \notin I$.
Det er her tilstrækkeligt at vælge $\epsilon$, således at $\epsilon |\textbf{a}_i^T\textbf{d}| < \textbf{a}_i^T\textbf{x} - b_i$ for alle $i \notin I$.
Når $\epsilon$ er lille nok, haves det, at $\textbf{y} \in \mathcal{P}$, og med samme argumentation haves, at $\textbf{z} \in \mathcal{P}$.
Dette betyder, at $\textbf{x}$ ikke er et ekstremumspunkt, da $\textbf{x} = \frac{\textbf{y} + \textbf{z}}{2}$.
Ved kontraponering ses (a) $\Rightarrow$ (c).
\\\\
%
Lad $\textbf{x}$ være en basal mulig løsning, lad $I = \{i \text{  } | \text{  } \textbf{a}_i^T\textbf{x} = b_i \}$ og lad $\textbf{x}=\sum_{i \in I} \textbf{a}_i$.
Så haves, at
%
\begin{align*}
\textbf{c}^T\textbf{x} = \sum_{i\in I}\textbf{a}_i^T\textbf{x} = \sum_{i\in I}b_i.
\end{align*}
%
Yderligere haves, at $\textbf{a}_i^T\textbf{x} \geq b_i$ for ethvert $\textbf{x} \in \mathcal{P}$ og ethvert $i$, og 
%
\begin{align}\label{eq:ulig}
\textbf{c}^T\textbf{x} = \sum_{i\in I}\textbf{a}_i^T\textbf{x} \geq \sum_{i\in I}b_i.
\end{align}
%
Dette viser, at $\textbf{x}$ er en optimal løsning til minimumsproblemet $\textbf{c}^T\textbf{x}$ over $\mathcal{P}$.
Yderligere holder uligheden i ligning \ref{eq:ulig}, hvis og kun hvis $\textbf{a}_i^T\textbf{x} = b_i$ for alle $i \in I$.
Da $\textbf{x}$ er en basal mulig løsning, er der $n$ lineært uafhængige betingelser, som er bindende i $x$, og $x$ er en entydig løsning til ligningssystemet $\textbf{a}_i^T\textbf{x} = b_i$, $i \in I$.
Det følger derfor, at $\textbf{x}$ er entydigt minimum til $\textbf{c}^T\textbf{x}$ over $\mathcal{P}$ og dermed et hjørnepunkt til $\mathcal{P}$.
Hermed (c) $\Rightarrow$ (b).
%
\end{proof}\\
%
Da en vektor er en basal mulig løsning, hvis og kun hvis det er et ekstremumspunkt og da definitionen på et ekstremumspunkt ikke refererer til en specifik repræsentation af polyeden, kan det konkluderes, at basale mulige løsninger er uafhængige af den valgte repræsentation.
%
\begin{kor}{}{}
Der kan kun være et endeligt antal af basale løsninger og basale mulige løsninger, givet der er et endeligt antal af lineære ulighedsbetingelser.
\end{kor}
%
\begin{proof}
Betragt et system af $m$ lineære ulighedsbetingelser for vektoren $\textbf{x} \in \R^n$.
Ved enhver basal løsning er der $n$ lineære ulighedsbetingelser, som er bindende.
Da $n$ bindende lineære ulighedsbetingelser definerer et entydigt punkt, følger det, at en anden basal løsning må have en anden mængde af $n$ bindende lineære ulighedsbetingelser.
Dette betyder, at antallet af basale løsninger har en øvre begrænsning, som afhænger af mængden af forskellige måder, hvorpå $n$ betingelser kan udvælges af de $m$ betingelser.
\end{proof}
%
\subsection{Tilstødende mulige løsninger}
%
To forskellige basale løsninger til samme mængde af lineære betingelser i $\R^n$ siges at være tilstødende, hvis der findes $n-1$ lineært uafhængige betingelser, som er bindende ved begge. 
Hvis de to basale løsninger er basale mulige løsninger, så er det forbindende linjestykke en af løsningsmængdens kanter.

% 
% 2.3 Polyede på standardform
% ---------------------------------------------------------
	\section{Polyeder på standardform}
\label{afsnit:fisk}
%
Som beskrevet i afsnit \ref{sec:standard}, kan optimeringsproblemer opskrives på standardform.
Dette korresponderer med polyeder, der ligeledes kan opskrives på standardform: 
$\mathcal{P}=\{ \mathbf{x} \in \R^n  \mid  A\mathbf{x}=\mathbf{b}, \mathbf{x} \geq \mathbf{0} \}$, hvor er $A$ er en $m \times n$ matrix.
$\mathcal{P}$ er her et \textit{polyeder på standardform}.
I de fleste tilfælde er det en fordel at antage, at de $m$ rækker i $A$ er lineært uafhængige.
I sætning \ref{thm:julieerfantalastisk} vises endvidere, at antagelsen om, at rækkerne er uafhængige, kan foretages, da de svarer til overflødige betingelser, som kan udelades.
Som det ses i definition \ref{defn:basal}, skal der være $n$ aktive begrænsninger i spil for at finde en basal mulig løsning.
Såfremt $m \neq n$, skal der derfor vælges $n-m$ variable $\textbf{x}_i$ med henblik på, at sætte disse $\textbf{x}_i=0$, hvilket gør begrænsningerne $\textbf{x}_i \geq 0$ aktive.
%ovenstående 53 i bogen har behov for andre læser så vi kan diskutere hvad det betyder.
Det er dog ikke uvæsentligt, hvilke af disse variable der omdannes til $0$ hvilket belyses af sætning \ref{thm:polystd}.
%
\begin{thm}{}{polystd}
Ved begrænsningerne $A\textbf{x}=\mathbf{b}$ og $\textbf{x}\geq 0$.
Antag, at $A$ er en $m \times n$ matrix og har lineært uafhængige rækker.
Vektoren $\textbf{x} \in \R^n$ er en basal løsning, hvis og kun hvis $A\textbf{x}=\mathbf{b}$, og der eksisterer indekser $B(1),\ldots,B(m)$, hvorom der gælder følgende:
\begin{enumerate}[label=(\alph*)]
\item Søjlerne $\mathbf{A}_{B(1)},\ldots,\mathbf{A}_{B(m)}$ er lineært uafhængige.
\item Hvis $i \neq B(1),\ldots, B(m)$, så er $\textbf{x}_i=0$.
\end{enumerate}
\end{thm}
\begin{proof}
Betragt et $x \in \R^n$ og antag, at der eksisterer indekser $\mathbf{A}_{B(1)},\ldots,\mathbf{A}_{B(m)}$, der opfylder (a) og (b).
%
Det gælder således, at de aktive begrænsninger, at $\textbf{x}_i=0$, når $i\neq B(1),\ldots,B(m)$, samt at $A\textbf{x}=\mathbf{b}$.
Dette implicerer, da det forrige gælder,  at 
%
$$\sum_{i=1}^{m}\textbf{A}_{B(i)}\textbf{x}_{B(i)}=\sum_{i=1}^{n}\textbf{A}_ix_i=A\textbf{x}=\textbf{b}.$$
%
Da søjlerne $\textbf{A}_{B(i)}$ for $i=1,\ldots,m$ er lineært uafhængige, kan $\textbf{x}_{B(1)},\ldots,\textbf{x}_{B(m)}$ bestemmes entydigt. 
Altså har ligningssystemet skabt af de aktive begrænsninger en entydig løsning.
%x_{B(i)} fra "da søjlerne ..." 
Jævnfør sætning \ref{thm:betingi},
% Ved ikke om denne reference er rigtig xD
følger det, at der er $n$ aktive begrænsninger, hvoraf $\mathbf{x}$ er en basal løsning. 
\\\\
%
%
Antag nu, at $\mathbf{x}$ er en basal løsning. 
Det skal nu vises, at (a) og (b) da er opfyldt.
Lad $\textbf{x}_{B_1},\ldots,\textbf{x}_{B_k}$ være ikke-nul komponenter i $\textbf{x}$.
Eftersom $\mathbf{x}$ er en basal løsning, følger nu, at ligningssysemet givet ved de aktive begrænsninger $\textbf{x}_i=0$, når $i\neq B(1),\ldots,B(k)$, samt  $\sum_{i=1}^{n}\mathbf{A}_i\textbf{x}_i=\mathbf{b}$, har en entydig løsning. 
Det samme må derfor gøre sig gældende for $\sum_{i=1}^{k}\mathbf{A}_{B(i)}\textbf{x}_{B(i)}=\mathbf{b}$.
Det følger derfor, at søjlerne i $A_{B(1)},\ldots,A_{B(k)}$ er lineært uafhængige.
%
Hvis dette ikke var tilfældet, ville der findes løsninger til $\sum_{i=1}^{k}\mathbf{A}_{B(i)} \lambda_i=\mathbf{0}$ udover den trivielle, hvor $\lambda_i=0$ for $i=1,2,\ldots,k$, hvilket ville betyder, at løsningen $\mathbf{x}$ ikke er entydig. 
Dette er i modstrid til at denne er en basal løsning.
$\mathbf{A}_{B(1)},\ldots ,\mathbf{A}_{B(k)}$ er således lineært uafhængige og $k \leq m$.
Da $A$ har $m$ lineært uafhængige rækker, er der ligeledes $m$ lineært uafhængige søjler.
%her kommer noget der følger fra sætning 1.3 i sektion 1.5 tror ikke vi har noget tilsvarende.
Der kan derfor findes $m-k$ søjler $A_{B(k+1)},\ldots,A_{B(m)}$, hvorom det gælder, at søjlerne $\mathbf{A}_{B(i)}$ med $i=1,\ldots,m$, er lineært uafhængige.
Da $k \leq m$ gælder det, hvis $i \neq B(1),\ldots,B(m)$, at $i \neq B(1),\ldots,B(k)$ og $x_i=0$.
%
\end{proof}
%der skal flere igennem det her bevis, tænker det er svært for andre end mig.
\\
\noindent
Som følge af sætning \ref{thm:polystd} kan basale løsninger konstrueres ved hjælp af følgende procedure.
\begin{enumerate}
\item Vælg $m$ lineært uafhængige søjler $\textbf{A}_{B(1)},\ldots,\textbf{A}_{B(m)}.$
\item Lad $\textbf{x}_i=0$ for $i \neq B(1),\ldots,B(m).$
\item Løs ligningssystemet med de $m$ ligninger, $A\textbf{x}=\textbf{b}$ for $\textbf{x}_{B(1)}, \ldots , \textbf{x}_{B(m)}.$
\end{enumerate}
Hvis løsningen, der er konstrueret af denne procedure, er ikke-negativ, så er løsningen en basal mulig løsning. Søjlerne $\textbf{A}_{B(1)},\ldots,\textbf{A}_{B(m)}$ kaldes for basis variablerne. Hvert valg af de lineært uafhængige søjler svarer til forskellige baser for rummet.
\\\\
\noindent
Det bevises nu, at antagelsen om at rækkerne var lineært uafhængige, kunne fortages.
\begin{thm}{}{julieerfantalastisk}
Lad $P=\{\textbf{x} \mid  A\textbf{x}=\textbf{b},x \geq 0\}$ være et ikke-tom polyeder, hvor $A$ er en $m \times n$ matrix med rækker $\textbf{a}^{T}_{1},\ldots,\textbf{a}^{T}_{m}$.
Antag, at rang$(A)=k<m$ og at rækkerne $\textbf{a}^T_{i_1},\ldots,\textbf{a}^T_{i_k}$ er lineært uafhængige. Givet polyederet 
$$Q=\{\textbf{x} \mid \textbf{a}^T_{i_1}\textbf{x}=b_{i_1},\ldots,\textbf{a}^T_{i_k}\textbf{x}=b_{i_k}, \textbf{x} \geq 0  \}$$ 
så er $Q=P$.
\end{thm}
\begin{proof}
Det bevises i tilfældet af $i_1=1,\ldots,i_k=k$, som er tilfældet, hvor de første $k$ rækker er lineært uafhængige. 
Alle andre tilfælde kan omskrives til dette tilfælde ved rækkeombytning. 
Det gælder, at $\mathcal{P}$ er en delmængde af $Q$, da alle elementer i $\mathcal{P}$ opfylder $Q$'s betingelser. 
Det skal dernæst vises, at $Q$ er en delmængde af $\mathcal{P}$.
Da rang$(A)=k$, har rækkerummet af $A$ dimension $k$ og et basis bestående af søjlerne $\textbf{a}_1,\ldots,\textbf{a}_k$. 
Derfor kan kan alle rækker $\textbf{a}_i$ af $A$ blive udtrykt som linearkombination af de andre rækker ved $\textbf{a}_i=\sum^{k}_{j=1}\lambda_{ij}\textbf{a}_j$ for skalarer $\lambda_{ij}$. 
Lad $\textbf{x}$ være et element i $\mathcal{P}$. 
%Hvilke(t) x skal være med fed her?
Så gælder
$$b_i=\textbf{a}_i\textbf{x}=\sum^{k}_{j=1}\lambda_{ij}\textbf{a}_j\textbf{x}=\sum^{k}_{j=1}\lambda_{ij}b_j,$$
for $i=1,\ldots,m.$
Betragt nu et element $y$ i $Q$. 
Dette vil også være tilhøre $\mathcal{P}$, da
$$ \textbf{a}_i\textbf{y}=\sum^{k}_{j=1}\lambda_{ij}\textbf{a}_j\textbf{y}=\sum^{k}_{j=1}\lambda_{ij}b_j=b_i.$$
Dermed er det vist, at $Q$ er en delmængde til $\mathcal{P}$. 
Dermed vist at $Q=\mathcal{P}$, da $\mathcal{P}$ ligeledes er en delmængde af $Q$.
\end{proof} \\	

% SIMPLEX METODEN
	\chapter{Simplex metoden}
% God arbejdslyst 	

% LU - DEKOMPOSITION 
	\chapter{LU-dekomposition}
% God arbejdslyst 

% Appendicer indsættes inde i en appendices-blok og bliver nummereret med
% bogstaver i stedet for tal
\begin{appendices}
\chapter{Entydigheden af en matrix på reduceret trappeform}

\begin{thm}{}{redutrap}
Lad $A$ være en matrix med den reducerede trappeform $A_R$.
Så er følgende sandt:
\begin{enumerate}[label=(\alph*)]
\item  Hvis en søjle $j$ i $A$ er en linearkombination  af andre søjler i $A$, så er søjlen i $A_R$ en linearkombination af de tilsvarende søjler i $A_R$ med de samme koefficienter.
\item  Hvis en søjle $j$ i $A_R$ er en linearkombination  af andre søjler i $A_R$, så er søjlen i $A$ en linearkombination af de tilsvarende søjler i $A$ med de samme koefficienter.
\end{enumerate}
\end{thm}%
\begin{proof}
Jævnfør \ref{thm:eleinv} så gælder det, at der findes en invertibel matrix $P$, så $PA=A_R$.
Deraf gælder $P\textbf{a}_i=\textbf{r}_i$ for alle $i$.
Antag, at søjlen $j$ i $A$ er en linearkombination af de andre søjler i $A$.
Så findes skalarer $c_1,c_2,\ldots,c_k$, så det gælder, at \\
\begin{equation}
\textbf{a}_j=c_1\textbf{a}_1+c_2\textbf{a}_2+\cdots+c_k\textbf{a}_k.
\end{equation}
Derfor 
\begin{align*}
\textbf{r}_j=P\textbf{a}_j&=P(c_1\textbf{a}_1+c_2\textbf{a}_2+\cdots+c_k\textbf{a}_k) \\
&= c_1P\textbf{a}_1+c_2P\textbf{a}_2+\cdots+c_kP\textbf{a}_k\\
&= c_1\textbf{r}_1+c_2\textbf{r}_2+\cdots+c_k\textbf{r}_k.
\end{align*}
Beviset for (b) gør brug af sammen fremgangsmåde, men benytter dog $\textbf{r}_i=P^{-1}\textbf{a}_i.$
\end{proof}\\
%
Heraf kan der udledes en række egenskaber for en matrix på reduceret trappeform.
%
\begin{lem}{}{denderlort}
Lad $A_R$ være en $m \times n$ matrix på reduceret trappeform.
Så gælder, at
\begin{enumerate}[label=(\alph*)]
\item En søjle i $A_R$ er en pivotsøjle, hvis og kun hvis den er ikke-nul og ikke er en linearkombination af de foregående søjler i $A_R$. 
\item Den $j$'te pivotsøjle i $A_R$ er $\textbf{e}_j \in \R^m$ og derfor er pivotsøjlerne i $A_R$ lineært uafhængige. 
\item Antag, at $\textbf{r}_j$ ikke er en pivotsøjle i $A_R$, og at der er $k$ pivotsøjler i $A_R$ før $\textbf{r}_j$.
Så gælder det, at $\textbf{r}_j$ er en linearkombination af de $k$ forudgående pivotsøjler og koefficienterne er de $k$ første indgange i $\textbf{r}_j$. 
\end{enumerate}
\end{lem}
\noindent
%
\ref{lem:denderlort} følger direkte af \ref{defn:trap} og \ref{defn:redtrap}
og kan generaliseres til \ref{lem:pivotu}, som bruges til \ref{thm:entydig}.
%
\begin{lem}{}{pivotu}
Følgende udsagn er sande for enhver matrix $A$: 
\begin{enumerate}[label=(\alph*)]
\item Pivotsøjlerne i $A$ er lineært uafhængige. 
\item Enhver ikke-pivotsøjle i $A$ er en linearkombination af de forrige pivotsøjler i $A$, hvor koefficienterne i linearkombinationen er indgangene i den tilsvarende søjle i den reducerede trappeform af $A$. 
\end{enumerate} 
\end{lem}
%
\begin{thm}{}{entydig}
Enhver matrix kan, ved hjælp af elementære rækkeoperationer, kun blive omdannet til en entydig matrix på reduceret trappeform.
\end{thm}
%
\begin{proof}
I følgende bevis refereres der til de tidligere egenskaber for en matrix på reduceret trappeform.\\\\
Lad $A$ være en matrix og $A_R$ være den reducerede trappeform af $A$. 
Jævnfør \ref{lem:denderlort}(a) gælder det, at en søjle i $A_R$ er en pivotsøjle i $A_R$, hvis og kun hvis den er en ikke-nulsøjle og ikke er en linearkombination af de forudgående pivotsøjler i $A_R$. 
Jævnfør \ref{thm:redutrap} gælder det, at hvis en søjle i $A_R$ er en pivotsøjle, så er den ikke en linearkombination af de forudgående søjler i $A_R$ og deraf er søjlen i $A$ heller ikke en linearkombination af de forudgående  søjler i $A$, hvilket betyder at positionerne af pivotsøjlerne i $A_R$ er entydigt bestemt af søjlerne i $A$. 
Endvidere gælder det, jævnfør \ref{lem:denderlort}(b), at den $j$'te pivotsøjle i $A_R$ er den $j$'te standardvektor i $\R^m$, og derfor er pivotsøjlerne i $A_R$ endvidere entydigt bestemt af søjlerne i $A$. \\\\
Dernæst bevises, at ikke-pivotsøjlerne i $A_R$ også er bestemt af søjlerne i $A$. 
Antag, at $\textbf{r}_j$ ikke er en pivotsøjle i $A_R$.
Hvis $\textbf{r}_j = \textbf{0}$ medfører det, at $\textbf{a}_j = \textbf{0}$ jævnfør \ref{thm:redutrap}. 
Antag nu, at $\textbf{r}_j \neq \textbf{0}$.
Så gælder det jævnfør \ref{lem:denderlort}(c), at $\textbf{r}_j$ er en linearkombination af de forudgående pivotsøjler i $A_R$, som er lineært uafhængige. 
Det gælder endvidere, at koefficienterne i linearkombinationen er indgangene i $\textbf{r}_j$, hvor indgangene, som ville være rækkens første ikke-nul-indgang, er nulindgange.
Jævnfør \ref{thm:redutrap} er $\textbf{a}_j$ ligeledes en linearkombination af de forudgående pivotsøjler i $A$, som er lineært uafhængige og har de samme koefficienter, der afhænger af $A$. 
Derfor er $\textbf{r}_j$ entydigt bestemt af $A$ og det konkluderes deraf, at $A_R$ er entydig bestemt af $A$. 
\end{proof}
\chapter{vigtige egenskaber}

Følgende udvalgte egenskaber ved lineært afhængige og uafhængige mængder, er vigtige for visse beviser senere i rapporten.
\subsubsection{Egenskaber af lineært afhængige og uafhængige mængder}\label{egenskab_lin}
\begin{enumerate}
\item En mængde bestående af én ikke-nul er lineært uafhængig, men $\{\textbf{0}\}$ er lineært afhængig.
\item Ethvert underrum af $\R^n$ der indeholder mere end $n$ vektorer er lineært afhængig
\item Hvis spannet $\S$ er et underrum af $\R^n$ og ingen vektorer kan fjernes uden at ændre spannet, så er $\S$ lineært uafhængig

\end{enumerate}

\end{appendices}

% Dokumentets 'back matter' er til ekstra ting som f.eks. litteraturlisten.
% Overskrifter bliver ikke nummereret her.
\backmatter

% Automatisk litteraturliste baseret på, hvilke kilder, der er blevet refereret
% til i løbet af rapporten.
\bibliographystyle{apalike}
\bibliography{
  incl/bib/books,
  incl/bib/articles,
  incl/bib/software
}

\end{document}
